% !TeX root = ../../Book1.tex
\chapter{迹理论与分数阶 Sobolev 空间}
\label{b1:ch:23}
\BookChapterNumber{b1:ch:23}
\begin{chapterabstract}
本章从一个边界值问题出发：域内几乎处处等价的 Sobolev 函数，怎样决定边界上的数据？先通过沿线端点与范数估计构造迹算子，证明有限指数下零迹与紧支集光滑闭包等价；再以点对差分定义分数阶 Sobolev 空间，刻画一阶迹的真实目标，并构造有界线性右逆。最后讨论平坦边界上的高阶数据相容性、法向导数在折角处的限制，以及 Besov 空间与本章差分范数的关系。
\end{chapterabstract}

\begin{learninggoals}
\begin{enumerate}
\item\label{b1:goal:23-trace}区分可测代表的边界赋值与由内部正则性确定的迹，完成半空间和有界 Lipschitz 区域上的 \(L^p\) 迹构造。
\item\label{b1:goal:23-zero}通过局部压平、零端点拼合及内部逼近，证明有限指数下的零迹刻画，说明其中的几何和指数条件。
\item\label{b1:goal:23-fractional}掌握分数阶双积分、差分表示、完备性、乘子、坐标变换及光滑稠密性，追踪非局部截断的额外项。
\item\label{b1:goal:23-boundary-space}证明 \(1<p<\infty\) 时的一阶分数迹定理与满射性，理解右逆的尺度卷积构造及边界范数的图表独立性。
\item\label{b1:goal:23-higher}在选读部分辨认高阶边界数据之间的切向相容关系，区分平坦法向和变化的法向场，并理解当前分数阶空间的 Besov 表述。
\end{enumerate}
\end{learninggoals}

本章使用\chapref{b1:ch:21}的沿线绝对连续代表，以及\chapref{b1:ch:22}的逼近、补零和一阶延拓。边界面积通过 Lipschitz 图像公式计算。因此先前的几何测度工具并未退出分析问题：它们保证不同边界参数中所说的“几乎处处”相容，也使边界积分得到正确的面积因子。

初读可先读完\secref{b1:sec:2301}，建立零迹与 \(W_0^{1,p}\) 的联系；再依次读分数阶空间的基本工具和完整的一阶迹定理。\secref{b1:sec:2304}为选读，后续内部嵌入理论不依赖其中未展开的高阶右逆或一般 Besov 理论。特别应留意两个范围：本章的 \(L^p\) 迹与零迹刻画包含 \(p=1\)，精确分数阶目标及所构造的线性右逆则只在 \(1<p<\infty\) 证明。

中文伴读可选王明新的《索伯列夫空间》\cite{Wang2013SobolevChinese}，先对照第2.4节、第2.10节的边界迹，再读第3.5节中的实指数空间与迹。该书书名采用人名的中文音译，本书正文仍沿用 Sobolev 原拼。Di Nezza、Palatucci 与 Valdinoci 的《Hitchhiker's guide to the fractional Sobolev spaces》\cite{DiNezza2012Hitchhikers}适合伴读本章的差分部分；涉及原文仅陈述或采用不同前置工具的地方，各节另行说明，不把一条参考文献当作未写出的证明。

% !TeX root = ../../Book1.tex
\section{Sobolev 迹}
\label{b1:sec:2301}

Sobolev 函数首先是域内的几乎处处等价类。给这个等价类随意补上边界值，不会改变任何体积积分，却会任意改变所谓“边界限制”。所以迹不是把某个可测代表直接放到边界上，而是一个由内部正则性决定、并对 Sobolev 范数连续的运算。

本节先把目标空间取为边界上的 \(L^p\)，并处理全部有限指数 \(1\le p<\infty\)。更精细的分数阶目标留到\secref{b1:sec:2303}。暂时不要求读者预先理解分数阶空间，也不把有限指数的零迹刻画推广到 \(p=\infty\)。

\subsection{光滑函数的边界限制}
\label{b1:subsec:230101}

设 \(\Omega\subset\mathbb R^d\) 为有界 Lipschitz 区域，记
\[
 \sigma(A)=\mathcal H^{d-1}(A\cap\partial\Omega).
\]
这里表示将 \((d-1)\) 维 Hausdorff 测度限制到边界；以后直接写 \(\int_{\partial\Omega}f\,\mathrm d\sigma\)。在\secref{b1:sec:2203}的图表中，边界参数为
\[
 \Gamma(z)=(z,\varphi(z)).
\]
由 Lipschitz 图像的面积公式，
\begin{equation}\label{b1:eq:lipschitz-boundary-measure}
 \int_{\Gamma(A)}f\,\mathrm d\sigma
 =\int_A f(\Gamma(z))\sqrt{1+|\nabla\varphi(z)|^2}\,\mathrm dz.
\end{equation}
梯度在几乎处处存在，面积因子处于 \(1\) 与 \(\sqrt{1+L^2}\) 之间。因此图表中的零测集、可积性和 \(L^p\) 范数都能与横向 Lebesgue 测度比较。有限图表覆盖又保证 \(\sigma(\partial\Omega)<\infty\)。维数为一时，\(\mathcal H^0\) 是计数测度，有界 Lipschitz 区域的边界是有限点集。

若 \(v\) 在 \(\overline\Omega\) 的一个邻域内光滑，则 \(v|_{\partial\Omega}\) 没有歧义。把它延伸到一般 Sobolev 元素之前，需要先知道这样的函数确实足够多。

\begin{proposition}\label{b1:prop:sobolev-smooth-up-to-boundary}
设 \(1\le p<\infty\)，且 \(\Omega\) 是有界 Lipschitz 区域。则
\[
 \{\,v|_\Omega\mid v\in C_c^\infty(\mathbb R^d)\,\}
\]
在 \(W^{1,p}(\Omega)\) 中稠密。
\end{proposition}
\begin{proof}
取\cref{b1:thm:lipschitz-domain-extension}的有界延拓 \(Eu\in W^{1,p}(\mathbb R^d)\)。由全空间光滑稠密性，选 \(v_j\in C_c^\infty(\mathbb R^d)\)，使 \(v_j\to Eu\) 于 \(W^{1,p}(\mathbb R^d)\)。限制到 \(\Omega\) 不增加范数，且 \((Eu)|_\Omega=u\)，所以 \(v_j|_\Omega\to u\)。
\end{proof}

这里的紧支集属于整个空间，不必留在 \(\Omega\) 内；这些近似函数可以在边界处非零。这与 \(W_0^{1,p}\) 的定义有本质区别。稠密性本身还不足以保证它们的边界限制收敛，必须另外建立边界范数估计。

\subsection{迹算子}
\label{b1:subsec:230102}

先在 \(H_+=\mathbb R^{d-1}\times(0,\infty)\) 上说明边界估计的来源。横向变量写作 \(x'\)，纵向变量写作 \(t\)。沿几乎每条竖线，Sobolev 函数具有绝对连续代表，因而有机会通过 \(t\downarrow0\) 确定边界值。

\begin{theorem}[半空间的 \(L^p\) 迹]\label{b1:thm:half-space-lp-trace}
设 \(1\le p<\infty\)。存在唯一有界线性算子
\[
 \operatorname{Tr}_+:W^{1,p}(H_+)\longrightarrow L^p(\mathbb R^{d-1}),
\]
使它对 \(C_c^\infty(\mathbb R^d)\) 的限制等于 \(u(x',0)\)。它由几乎每条竖线的绝对连续端点值给出，并满足
\begin{equation}\label{b1:eq:half-space-lp-trace}
 \|\operatorname{Tr}_+u\|_{L^p(\mathbb R^{d-1})}
 \le\|u\|_{L^p(\mathbb R^{d-1}\times(0,1))}
      +\|\partial_du\|_{L^p(\mathbb R^{d-1}\times(0,1))}.
\end{equation}
选择这些沿线代表后，对每个 \(0<t<1\)，
\begin{equation}\label{b1:eq:trace-slice-convergence}
 \|u(\cdot,t)-\operatorname{Tr}_+u\|_{L^p}
 \le t^{1-1/p}
       \|\partial_du\|_{L^p(\mathbb R^{d-1}\times(0,t))}.
\end{equation}
特别地，切片在 \(L^p\) 中趋于迹；\(p=1\) 时右侧也趋零。
\end{theorem}
\begin{proof}
取\cref{b1:thm:sobolev-acl}的共同代表。Fubini 定理保证，对几乎每个 \(x'\)，函数 \(u(x',\cdot)\) 与 \(\partial_du(x',\cdot)\) 都在 \((0,1)\) 上可积。因此其绝对连续代表延伸到端点，并可写成
\[
 u(x',t)=a(x')+\int_0^t\partial_du(x',s)\,\mathrm ds.
\]
在例外的 \(x'\) 上赋零；端点值可由 \(u(x',1/j)\) 的极限取得，故可取为可测函数。由上式，对 \(0<t<1\)，
\[
 |a(x')|\le |u(x',t)|
                  +\int_0^1|\partial_du(x',s)|\,\mathrm ds.
\]
对 \(t\in(0,1)\) 平均，再对 \(x'\) 取 \(L^p\) 范数。积分 Minkowski 不等式和区间长度为一给出\eqref{b1:eq:half-space-lp-trace}。因此 \(a\in L^p\)，令 \(\operatorname{Tr}_+u=a\)。两代表几乎处处相同，则对几乎每个 \(x'\) 沿线几乎处处相同，绝对连续代表及其端点值因而一致。这使定义在等价类层面良定义并线性。

在端点积分表示中用 Hölder 不等式，
\[
 |u(x',t)-a(x')|^p
 \le t^{p-1}\int_0^t|\partial_du(x',s)|^p\,\mathrm ds.
\]
对 \(x'\) 积分即得\eqref{b1:eq:trace-slice-convergence}。当 \(p=1\) 时，用积分三角不等式取得同一结论，右端由可积函数在缩小边界层上的积分趋零。沿线代表的选择还使这些切片对每个 \(t\) 都可测并属于 \(L^p\)。

对光滑函数，端点值当然就是通常限制。最后，将 \(u\) 偶反射到全空间，再用全空间 \(C_c^\infty\) 稠密性，说明光滑限制在 \(W^{1,p}(H_+)\) 中稠密。任意两个有界线性算子若在这一稠密类上相同，就在全空间相同，得到唯一性。
\end{proof}

这种由连续性延伸边界限制的算子，称为\term[ji-suanzi]{迹算子}[trace operator]，其输出称为函数的\term[ji]{迹}[trace]，常记为 \(\operatorname{Tr}u\) 或 \(\gamma_0u\)。本书统一以前一记号表示零阶边界数据。“迹”不是函数在任意可测代表上的边界赋值；上述定理通过内部沿线积分把它唯一地确定下来。
\symidx{\(\operatorname{Tr}u\)：Sobolev 函数的边界迹}

\subsection{\texorpdfstring{\(W^{1,p}\)}{W1,p} 的迹}
\label{b1:subsec:230103}

对于 Lipschitz 边界，竖线改成图表中横向位置固定的线段。图函数只需 Lipschitz，因为本章首先处理一阶空间，双 Lipschitz 换元已有相应弱导数规则。

\begin{theorem}[Lipschitz 区域的 \(L^p\) 迹]\label{b1:thm:lipschitz-lp-trace}
设 \(\Omega\subset\mathbb R^d\) 为有界 Lipschitz 区域，\(1\le p<\infty\)。存在唯一有界线性算子
\[
 \operatorname{Tr}:W^{1,p}(\Omega)\longrightarrow L^p(\partial\Omega,\sigma)
\]
延伸光滑函数的通常边界限制，并满足
\[
 \|\operatorname{Tr}u\|_{L^p(\partial\Omega)}
 \le C_{\Omega,p}\|u\|_{W^{1,p}(\Omega)}.
\]
若 \(u\in W^{1,p}(\Omega)\cap C(\overline\Omega)\)，则迹与该连续代表的边界限制相同。在各个局部图表内，迹也等于相应竖线的绝对连续端点值。
\end{theorem}
\begin{proof}
先设 \(d\ge2\)。取\secref{b1:sec:2203}的有限内外图表、压平映射 \(F_j\) 及光滑分割函数 \(\eta_j\)。内部项 \(\eta_0\) 在边界附近为零，其他项满足 \(\sum_{j=1}^N\eta_j=1\) 于边界。记 \(\Gamma_j(z)=F_j(z,0)\)。

对 \(u\in W^{1,p}(\Omega)\)，先将 \(\eta_ju\) 拉回正半盒。它在侧面及顶面附近为零，可以像局部延拓引理那样补零为 \(H_+\) 上的 \(W^{1,p}\) 函数 \(v_j\)。这里不跨过 \(t=0\) 补零，故不需要预先假设零边界值。定义
\[
 T_ju=(\operatorname{Tr}_+v_j)\circ\Gamma_j^{-1}
\]
于该边界片，并在其余边界上取零。由半空间迹估计、坐标换元、光滑乘法及\eqref{b1:eq:lipschitz-boundary-measure}，
\[
 \|T_ju\|_{L^p(\partial\Omega)}
 \le C_j\|u\|_{W^{1,p}(\Omega)}.
\]
有限和 \(Tu=\sum_{j=1}^N T_ju\) 因而是有界线性映射。

若 \(u\) 光滑到边界，则沿图表竖线有通常连续端点值，即
\[
 T_ju=(\eta_j|_{\partial\Omega})(u|_{\partial\Omega}).
\]
分割函数之和在边界为一，所以 \(Tu=u|_{\partial\Omega}\)。由\cref{b1:prop:sobolev-smooth-up-to-boundary}的稠密性，\(T\) 是唯一延伸光滑限制的有界算子，记为 \(\operatorname{Tr}\)。不同图表或分割函数得到的构造必然一致。

还需说明局部相容性并不限于本次选定的分割函数。给定一个图表及支集留在其内的光滑截断 \(\chi\)，取全空间光滑限制 \(u_m\to u\) 于 \(W^{1,p}(\Omega)\)。弱乘积与双 Lipschitz 换元保证 \((\chi u_m)\circ F\to(\chi u)\circ F\) 于相应正半盒。对各 \(u_m\)，两种端点解释一致；两边的 \(L^p\) 迹映射都连续，因此取极限仍一致。在任意较小边界片选取 \(\chi=1\)，即可得到所述局部端点刻画。

若 \(u\) 还有闭域连续代表，则沿这些图表竖线的端点值就是它的连续边界值，故两者一致。最后，\(d=1\) 时区域是有限个有界区间之并；在每个区间使用\secref{b1:sec:2202}的端点连续映射，再对有限端点求和，给出同一结论。
\end{proof}

\begin{proposition}\label{b1:prop:trace-multiplication}
若 \(\chi\) 在 \(\overline\Omega\) 的一个邻域内光滑，则
\[
 \operatorname{Tr}(\chi u)
 =(\chi|_{\partial\Omega})\operatorname{Tr}u
 \quad\text{于 }L^p(\partial\Omega).
\]
\end{proposition}
\begin{proof}
对光滑限制该等式只是通常乘法。对一般 \(u\)，用光滑限制依 \(W^{1,p}\) 逼近。由于 \(\overline\Omega\) 紧，\(\chi\) 及其一阶导数有界，乘法在 \(W^{1,p}(\Omega)\) 中连续；边界乘法及迹也连续，故等式可以取极限。
\end{proof}

迹定理并不说每个 \(L^p\) 边界函数都是某个 \(W^{1,p}\) 函数的迹。在 \(1<p<\infty\) 时，梯度可积性还强迫边界数据具有分数阶正则性；后两节将刻画这一附加条件。本节只先建立一个足以比较边界值、且涵盖 \(p=1\) 的连续目标空间。

\subsection{零迹与 \texorpdfstring{\(W_0^{1,p}\)}{W01,p}}
\label{b1:subsec:230104}

现在可以把上一章的闭包条件与真正的边界数据联系起来。零迹条件比“在边界上给代表赋零”强得多：它要求由内部弱导数确定的端点值为零。

\begin{lemma}\label{b1:lem:half-space-zero-trace-extension}
对 \(u\in W^{1,p}(H_+)\)、\(1\le p<\infty\)，若 \(\operatorname{Tr}_+u=0\)，则零延拓 \(\widetilde u\) 属于 \(W^{1,p}(\mathbb R^d)\)，并有
\[
 \partial_i\widetilde u=\widetilde{\partial_i u}.
\]
此外，\(u\in W_0^{1,p}(H_+)\)。
\end{lemma}
\begin{proof}
沿几乎每条竖线，\(u\) 的绝对连续代表在零点取值零。把它在负半线补零，两侧绝对连续函数的端点值相同，因而在跨过零点的紧区间上仍绝对连续；导数就是原导数的补零。沿其他坐标方向，正半空间保持原来的沿线性质，负半空间恒为零，中间整层直线在横向参数中是零测集。因此零延拓具有一个共同 ACL 代表，函数和候选导数又都属于 \(L^p\)。由 ACL 判据得到所求弱导数等式。

为得到内部光滑逼近，令 \(F=\widetilde u\)，取向内平移 \(F_h(x)=F(x-he_d)\)。函数及其弱导数的 \(L^p\) 平移连续性给出 \(F_h\to F\) 于 \(W^{1,p}(\mathbb R^d)\)，而 \(F_h=0\) 于 \(\{x_d<h\}\)。固定 \(h\)，用远处截断 \(\chi_R\) 令 \(\chi_RF_h\to F_h\)；这与全空间稠密性证明中的尾部估计相同。再以小于 \(h/2\) 的尺度磨光，所得函数的紧支集留在 \(H_+\) 内，且在全空间 Sobolev 范数中逼近 \(\chi_RF_h\)。依次令平移、尾部及磨光误差趋零，再限制回半空间，即得 \(W_0^{1,p}\) 的闭包条件。
\end{proof}

\begin{theorem}[零迹刻画]\label{b1:thm:zero-trace-characterization}
设 \(\Omega\) 为有界 Lipschitz 区域，\(1\le p<\infty\)。则
\[
 W_0^{1,p}(\Omega)
 =\{\,u\in W^{1,p}(\Omega)\mid\operatorname{Tr}u=0\,\}
 =\ker\operatorname{Tr}.
\]
\end{theorem}
\begin{proof}
若 \(u\in W_0^{1,p}(\Omega)\)，取 \(\varphi_m\in C_c^\infty(\Omega)\) 在 \(W^{1,p}\) 中逼近它。各 \(\varphi_m\) 在边界附近为零，故其迹均为零。迹的连续性给出 \(\operatorname{Tr}u=0\)。

反过来，设 \(\operatorname{Tr}u=0\)。维数为一的情形已由区间零端点刻画证明；以下设 \(d\ge2\)。仍用有限分解 \(u=\eta_0u+\sum_{j=1}^N\eta_ju\)。内部项具有内部紧支集，故由有限指数的内部磨光属于 \(W_0^{1,p}(\Omega)\)。

对一个边界项，将 \(\eta_ju\) 拉回正半盒并在半空间内补零，记为 \(v_j\)。迹的乘法及图表相容性说明 \(\operatorname{Tr}_+v_j=0\)，所以前一引理给出 \(v_j\in W_0^{1,p}(H_+)\)。取 \(a_m\in C_c^\infty(H_+)\)，使 \(a_m\to v_j\) 于 \(W^{1,p}(H_+)\)。

为在原图表内推回，选 \(\zeta\in C_c^\infty(D_j)\)，使它在 \(v_j\) 的闭支集附近等于一；这里支集可以接触 \(t=0\)，但仍紧含于外盒 \(D_j\)。光滑乘法连续，故 \(\zeta a_m\to v_j\) 于半空间 Sobolev 范数。每个 \(\zeta a_m\) 的支集同时避开 \(t=0\) 与外盒边界。

将 \(\zeta a_m\) 经 \(F_j^{-1}\) 推回，并在图表外补零，得到 \(\Omega\) 内部紧支集的 \(W^{1,p}\) 函数 \(b_m\)。虽然 Lipschitz 坐标变换不保证它光滑，但内部紧支集逼近保证 \(b_m\in W_0^{1,p}(\Omega)\)。由双 Lipschitz 换元的范数估计，\(b_m\to\eta_ju\) 于 \(W^{1,p}(\Omega)\)。闭性于是给出 \(\eta_ju\in W_0^{1,p}(\Omega)\)。有限个边界项与内部项相加，完成证明。
\end{proof}

证明中不能省略“推回后再次作内部光滑逼近”：图函数通常并不光滑，拉回坐标中的光滑函数推回去未必光滑。真正保留下来的是一阶 Sobolev 范数和内部紧支集，这已足以使用上一章的逼近定理。

有限 \(p\) 的范围也不可省去。区间函数 \(u(x)=x(1-x)\) 在两个端点取零，但\secref{b1:sec:2202}已经证明它不属于本书按范数闭包定义的 \(W_0^{1,\infty}(0,1)\)。另一方面，任意开集未必有单侧 Lipschitz 图表，上一章去掉中点的区间反例也不能被当前定理覆盖。零迹与闭包条件等价，既使用了指数范围，也使用了边界几何。

% !TeX root = ../../Book1.tex
\section{分数阶 Sobolev 空间}
\label{b1:sec:2302}

整数阶 Sobolev 范数通过弱导数衡量函数的变化。分数阶空间则直接比较两点的函数值：差值除以距离的一定幂，再对所有点对积分。这样既不必先定义一个“半阶导数”，也能记录介于函数值与一阶导数之间的正则性。稍后的迹空间将使用这一尺度，因为限制到边界后，正则性通常不再是整数阶。

Di Nezza、Palatucci 与 Valdinoci 的《Hitchhiker's guide to the fractional Sobolev spaces》\cite[Section 2 and Lemmas 5.1--5.3]{DiNezza2012Hitchhikers}给出双积分定义、指数比较及局部化工具。本节固定 \(0<s<1\)、\(1\leq p<\infty\)，函数可以取实值或复值；\(\Omega\subseteq\mathbb R^d\) 为非空开集，\(d\geq1\)。所有积分均使用 Lebesgue 测度，不在双积分前加入依赖于 \(s\) 的归一化系数。

% 实读 DiNezza2012Hitchhikers 的作者 arXiv:1104.4345v3 副本：
% 内页5--10，式(2.1)--(2.4)、Proposition2.1/2.2及Theorem2.4；
% 内页33--36，Lemma5.1--5.3及其证明。页码为作者副本内部页码，非期刊页码。
% 2.4仅陈述稠密性并转引Adams，未将其视为已读取的完整证明。
% 本节稠密性用差分映射的Lp平移连续性与远处截断另行完整展开；
% 双Lipschitz换元用本书已证测度变换界，不借后文迹理论。

\subsection{Slobodeckij 半范数}
\label{b1:subsec:230201}

\begin{definition}\label{b1:def:slobodeckij-seminorm}
设 \(u:\Omega\rightarrow\mathbb K\) 可测。定义其\term[slobodeckij-banfanshu]{Slobodeckij 半范数}[Slobodeckij seminorm]为
\begin{equation}\label{b1:eq:slobodeckij-seminorm}
 [u]_{s,p;\Omega}
 \coloneqq\left(\int_\Omega\int_\Omega
       \frac{|u(x)-u(y)|^p}{|x-y|^{d+sp}}\,\mathrm dy\,\mathrm dx\right)^{1/p},
\end{equation}
允许取值 \(+\infty\)。对角线 \(x=y\) 上的被积函数约定为零。全空间的半范数简记为 \([u]_{s,p}\)。
\end{definition}
\symidx{\([u]_{s,p;\Omega}\)：Slobodeckij 半范数}

Di Nezza 等在所引文章的 Section 2 将同一双积分称为 Gagliardo semi-norm，亦即以 Gagliardo 命名的半范数。\index[foreign]{Gagliardo semi-norm}这里的两个名称指向同一个差分条件；本书保持上述主称及不附加归一化系数的约定。

改变 \(u\) 在一个零测集 \(N\) 上的值，只会改变 \((N\times\Omega)\cup(\Omega\times N)\) 上的被积函数；这个集合由 Tonelli 定理为零测集。因此半范数只依赖几乎处处等价类。它的齐次性来自绝对值，三角不等式则是对函数
\[
 (x,y)\longmapsto\frac{u(x)-u(y)}{|x-y|^{d/p+s}}
\]
应用 \(L^p(\Omega\times\Omega)\) 中的 Minkowski 不等式。若半范数为零，Fubini 定理给出某个 \(y\in\Omega\)，使 \(u(x)=u(y)\) 对几乎每个 \(x\) 成立；反过来，常数函数的半范数为零。

这个结论不需要 \(\Omega\) 连通，因为积分也比较不同连通分支上的点。例如取 \(\Omega=(-2,-1)\cup(1,2)\)，令 \(u\) 在左区间等于 \(1\)，在右区间等于 \(0\)。它的弱导数在 \(\Omega\) 内为零，但
\[
 [u]_{s,p;\Omega}^p
 =2\int_{-2}^{-1}\int_1^2\frac{1}{|x-y|^{1+sp}}\,\mathrm dy\,\mathrm dx>0.
\]
两个区间的距离为 \(2\)，所以该积分有限。分数阶半范数不仅记录每个小邻域中的变化，还记录相隔位置之间的差异。

\begin{proposition}\label{b1:prop:fractional-difference-formula}
对 \(u\in L^p(\mathbb R^d)\)，沿用平移约定
\[
 \tau_hu(x)\coloneqq u(x-h),\qquad
 \Delta_hu\coloneqq\tau_hu-u.
\]
也就是说，\(\Delta_hu(x)=u(x-h)-u(x)\)。则
\begin{equation}\label{b1:eq:fractional-difference-formula}
 [u]_{s,p}^p
 =\int_{\mathbb R^d}\frac{\|\Delta_hu\|_p^p}{|h|^{d+sp}}\,\mathrm dh,
\end{equation}
两端允许同时为无穷。
\end{proposition}
\symidx{\(\tau_hu\)：平移，\(\tau_hu(x)=u(x-h)\)}
\symidx{\(\Delta_hu\)：有限差分，\(\Delta_hu=\tau_hu-u\)}
\begin{proof}
双积分中的函数非负。对每个固定 \(x\) 作平移换元 \(y=x-h\)，再用 Tonelli 定理交换 \(x,h\) 的积分次序，内层积分便是 \(\|\Delta_hu\|_p^p\)。无需预先假定半范数有限。
\end{proof}

因此，分数阶控制可以理解为：以不同距离平移函数，把各个平移误差按距离加权后合在一起。单个 \(h\) 上误差趋零，只是 \(L^p\) 平移连续性；这里还要求这些误差满足一个关于全部小尺度的可积条件。

\subsection{\texorpdfstring{\(W^{s,p}\)}{Ws,p}，\texorpdfstring{\(0<s<1\)}{0<s<1}}
\label{b1:subsec:230202}

\begin{definition}\label{b1:def:fractional-sobolev}
设 \(u\in L^p(\Omega)\)。若 \([u]_{s,p;\Omega}<\infty\)，则称 \(u\) 属于\term[fenshujie-sobolev-kongjian]{分数阶 Sobolev 空间}[fractional Sobolev space] \(W^{s,p}(\Omega)\)。该空间按几乎处处相等取商，赋予范数
\[
 \|u\|_{W^{s,p}(\Omega)}
 \coloneqq\left(\|u\|_{L^p(\Omega)}^p+[u]_{s,p;\Omega}^p\right)^{1/p}.
\]
\end{definition}
\symidx{\(W^{s,p}(\Omega)\)：分数阶 Sobolev 空间}

零阶项保证正定性。在无穷测度域上，它也控制远处的函数值，不能单凭双积分代替。

\begin{theorem}\label{b1:thm:fractional-sobolev-complete}
上述 \(W^{s,p}(\Omega)\) 是 Banach 空间。
\end{theorem}
\begin{proof}
记 \(G_su(x,y)=\bigl(u(x)-u(y)\bigr)/|x-y|^{d/p+s}\)，对角线上补零。映射 \(u\mapsto(u,G_su)\) 将本空间等距嵌入
\(L^p(\Omega)\times L^p(\Omega\times\Omega)\)，乘积使用两个范数的 \(p\) 次和。

设 \((u_j)\) 在 Sobolev 范数中为 Cauchy 列。由 \(L^p\) 完备性，存在 \(u\in L^p(\Omega)\) 和 \(F\in L^p(\Omega\times\Omega)\)，使 \(u_j\to u\)、\(G_su_j\to F\) 于相应 \(L^p\) 空间。可选共同子列 \((u_{j_\ell})\)，使
\[
 \sum_\ell\left(\|u_{j_\ell}-u\|_p^p+
                       \|G_su_{j_\ell}-F\|_p^p\right)<\infty.
\]
Tonelli 定理说明两组函数值误差分别几乎处处趋于零。删去 \(u_{j_\ell}\) 不趋于 \(u\) 的零测集所形成的两条乘积带后，\(G_su_{j_\ell}(x,y)\to G_su(x,y)\) 对几乎每个点对成立。因此 \(F=G_su\) 几乎处处，\(u\in W^{s,p}(\Omega)\)，而最初的整列在该范数中趋于 \(u\)。
\end{proof}

下面的三个工具用于把这个空间放到边界图表上。它们均直接控制点对积分，不要求对 \(u\) 作任何求导。

\begin{proposition}\label{b1:prop:fractional-lipschitz-multiplier}
设 \(a:\Omega\rightarrow\mathbb K\) 满足 \(|a|\leq A\) 且 Lipschitz 常数不超过 \(L\)。则乘法 \(u\mapsto au\) 在 \(W^{s,p}(\Omega)\) 上有界，具体有
\[
 [au]_{s,p;\Omega}^p
 \leq2^{p-1}A^p[u]_{s,p;\Omega}^p
       +C(d,s,p)(L^p+A^p)\|u\|_p^p,
 \quad \|au\|_p\leq A\|u\|_p.
\]
\end{proposition}
\begin{proof}
分解
\[
 a(x)u(x)-a(y)u(y)
 =a(x)\bigl(u(x)-u(y)\bigr)+u(y)\bigl(a(x)-a(y)\bigr).
\]
使用 \(|b+c|^p\leq2^{p-1}(|b|^p+|c|^p)\)，第一部分给出所列半范数项。对第二部分，利用
\(\bigl|a(x)-a(y)\bigr|\leq\min\{L|x-y|,2A\}\)。对每个 \(y\)，内层核积分不超过
\[
 L^p\int_{|h|<1}|h|^{p-d-sp}\,\mathrm dh
 +(2A)^p\int_{|h|\geq1}|h|^{-d-sp}\,\mathrm dh
 =d\omega_d\left(\frac{L^p}{p(1-s)}+\frac{(2A)^p}{sp}\right).
\]
这里用了\eqref{b1:eq:radial-one-dimensional-integral}，两个积分恰因 \(0<s<1\) 而有限。再乘以 \(|u(y)|^p\) 并积分，得到结论。
\end{proof}

\begin{proposition}\label{b1:prop:fractional-bilipschitz-change}
设 \(U,V\subseteq\mathbb R^d\) 为开集，\(\Phi:U\rightarrow V\) 为双 Lipschitz 同胚，\(\Phi\) 及其逆的 Lipschitz 常数分别不超过 \(L,M\)。则复合 \(u\mapsto u\circ\Phi\) 给出 \(W^{s,p}(V)\) 与 \(W^{s,p}(U)\) 之间的有界线性同构。正向估计为
\[
 \|u\circ\Phi\|_{L^p(U)}^p\leq M^d\|u\|_{L^p(V)}^p,
 \quad
 [u\circ\Phi]_{s,p;U}^p\leq L^{d+sp}M^{2d}[u]_{s,p;V}^p.
\]
\end{proposition}
\begin{proof}
逆映射保持 Lebesgue 零测集，故复合不依赖代表。由\eqref{b1:eq:sobolev-change-measure-bound}，对任意非负可测 \(b\)，有 \(\int_U b\bigl(\Phi(x)\bigr)\,\mathrm dx\leq M^d\int_Vb\)，从而得到零阶估计。

距离不等式 \(|\Phi(x)-\Phi(y)|\leq L|x-y|\) 给出
\[
 \frac{\bigl|u\bigl(\Phi(x)\bigr)-u\bigl(\Phi(y)\bigr)\bigr|^p}{|x-y|^{d+sp}}
 \leq L^{d+sp}
       \frac{\bigl|u\bigl(\Phi(x)\bigr)-u\bigl(\Phi(y)\bigr)\bigr|^p}{|\Phi(x)-\Phi(y)|^{d+sp}}.
\]
对右侧先在 \(x\) 变量使用测度变换界，再在 \(y\) 变量使用一次；Tonelli 定理允许这些非负积分运算，得到因子 \(M^{2d}\) 及所列半范数估计。对 \(\Phi^{-1}\) 重复论证，便得到逆方向的有界性。
\end{proof}

\begin{lemma}\label{b1:lem:compact-fractional-zero-extension}
设 \(u\in W^{s,p}(\Omega)\)，在 \(\Omega\setminus K\) 上几乎处处为零，其中 \(K\Subset\Omega\) 非空。记 \(\delta=\operatorname{dist}(K,\Omega^c)>0\)，在域外补零得到 \(\widetilde u\)。则 \(\widetilde u\in W^{s,p}(\mathbb R^d)\)，且
\begin{equation}\label{b1:eq:compact-fractional-zero-extension}
 [\widetilde u]_{s,p}^p
 \leq[u]_{s,p;\Omega}^p
       +\frac{2d\omega_d}{sp}\delta^{-sp}\|u\|_p^p,
 \quad\|\widetilde u\|_p=\|u\|_p.
\end{equation}
若 \(\Omega=\mathbb R^d\)，约定 \(\delta=\infty\)，附加项为零。
\end{lemma}
\begin{proof}
把全空间点对分成域内、域外及两个交叉部分，得到
\[
 [\widetilde u]_{s,p}^p
 =[u]_{s,p;\Omega}^p
   +2\int_K |u(x)|^p
      \left(\int_{\Omega^c}\frac{\mathrm dy}{|x-y|^{d+sp}}\right)\mathrm dx.
\]
若 \(\Omega^c\ne\varnothing\)，对 \(x\in K\) 和 \(y\in\Omega^c\) 有 \(|x-y|\geq\delta\)，故括号内不超过
\(\int_{|h|\geq\delta}|h|^{-d-sp}\,\mathrm dh
=d\omega_d\delta^{-sp}/(sp)\)。由此即得结论。空补集情形没有交叉积分，零阶范数等式则直接来自补零的定义。
\end{proof}

与整数阶内部补零不同，这里一般会增加半范数。原因是新加入的域外零值仍要与域内非零值比较；支集到边界的距离为这种相互作用提供了统一控制。边界图表中的截断必须在图表侧边留出余量，正是为了能够使用这一估计。

\subsection{基本嵌入}
\label{b1:subsec:230203}

\begin{proposition}\label{b1:prop:integer-to-fractional-sobolev}
全空间上有连续包含 \(W^{1,p}(\mathbb R^d)\subseteq W^{s,p}(\mathbb R^d)\)。具体地，
\begin{equation}\label{b1:eq:integer-fractional-estimate}
 [u]_{s,p}^p
 \leq d\omega_d\left(
    \frac{\|\nabla u\|_p^p}{p(1-s)}
    +\frac{2^p\|u\|_p^p}{sp}\right),
\end{equation}
其中 \(\|\nabla u\|_p\) 使用梯度的欧氏长度，与各偏导数 \(L^p\) 范数的有限乘积范数等价。
\end{proposition}
\begin{proof}
先设 \(u\in C_c^\infty(\mathbb R^d)\)。线段上的积分公式及 Hölder 不等式给出
\[
 \begin{aligned}
 u(x-h)-u(x)&=-\int_0^1\nabla u(x-th)\cdot h\,\mathrm dt,\\
 |u(x-h)-u(x)|^p
 &\leq |h|^p\int_0^1|\nabla u(x-th)|^p\,\mathrm dt.
 \end{aligned}
\]
对 \(x\) 积分，得到 \(\|\Delta_hu\|_p\leq|h|\cdot\|\nabla u\|_p\)。利用\cref{b1:thm:whole-space-sobolev-density}取 \(W^{1,p}\) 中的光滑逼近，函数差及梯度都在 \(L^p\) 中收敛，因而对每个固定 \(h\) 可把这一不等式传给一般 \(u\)。另一方面，平移不改变 \(L^p\) 范数，所以
\[
 \|\Delta_hu\|_p
 \leq\min\{\,|h|\cdot\|\nabla u\|_p,\ 2\|u\|_p\,\}.
\]
把它代入\eqref{b1:eq:fractional-difference-formula}。在 \(|h|<1\) 上用梯度界，在 \(|h|\geq1\) 上用函数值界；相应径向积分分别为 \(1/[p(1-s)]\) 和 \(1/(sp)\)，得到所列估计。
\end{proof}

这一论证先在全空间中进行，线段和平移都不会离开定义域。对已有有界 \(W^{1,p}\) 延拓算子的域，可以先延拓、再使用该命题并限制回原域；不能在任意开集上未经说明就沿两点间的整条线段积分。

估计中的两个分母说明了指数范围的作用：\(s\uparrow1\) 时，小尺度的径向积分接近 \(\int_0^1r^{-1}\,\mathrm dr\)；\(s\downarrow0\) 时，远处的径向积分接近 \(\int_1^\infty r^{-1}\,\mathrm dr\)。所以不能把这两个端点直接代入当前定义并认作原来的整数阶范数。适当归一化后的极限是另一项需要证明的结论。

\begin{proposition}\label{b1:prop:fractional-order-embedding}
若 \(0<t<s<1\)，则对任意开集 \(\Omega\)，有连续包含 \(W^{s,p}(\Omega)\subseteq W^{t,p}(\Omega)\)，并满足
\[
 [u]_{t,p;\Omega}^p
 \leq[u]_{s,p;\Omega}^p+\frac{2^pd\omega_d}{tp}\|u\|_p^p.
\]
若 \(\Omega\) 有界，另有更直接的估计
\( [u]_{t,p;\Omega}\leq(\operatorname{diam}\Omega)^{s-t}[u]_{s,p;\Omega}\)。
\end{proposition}
\begin{proof}
在 \(|x-y|<1\) 上，\(|x-y|^{-d-tp}\leq|x-y|^{-d-sp}\)，所以这一部分不超过 \(s\) 阶半范数的 \(p\) 次方。在余下部分，用
\( |u(x)-u(y)|^p\leq2^{p-1}(|u(x)|^p+|u(y)|^p)\)，并把每个核积分扩大到 \(|h|\geq1\)，便得附加项 \(2^pd\omega_d\|u\|_p^p/(tp)\)。如果直径有限，则直接在全部点对上使用 \(|x-y|^{(s-t)p}\leq(\operatorname{diam}\Omega)^{(s-t)p}\)，得到第二个估计。
\end{proof}

\subsection{光滑函数的稠密性}
\label{b1:subsec:230204}

对每个 \(h\)，普通 \(L^p\) 平移连续性已经成立，但若要在带奇异权的所有 \(h\) 上积分，还需要统一处理这些差分。把它们放进一个乘积空间即可完成这一点。

\begin{proposition}\label{b1:prop:fractional-translation-mollification}
设 \(u\in W^{s,p}(\mathbb R^d)\)，对上述平移有
\[
 \|\tau_zu-u\|_{W^{s,p}(\mathbb R^d)}\longrightarrow0\quad(z\to0).
\]
对\secref{b1:sec:2003}的磨光函数，还成立
\[
 u_\varepsilon=\rho_\varepsilon*u\in C^\infty(\mathbb R^d)\cap W^{s,p}(\mathbb R^d),
 \quad\|u_\varepsilon\|_{W^{s,p}}\leq\|u\|_{W^{s,p}},
 \quad u_\varepsilon\to u\text{ 于 }W^{s,p}.
\]
\end{proposition}
\begin{proof}
在 \(h\ne0\) 时定义
\[
 D_su(x,h)=\frac{\Delta_hu(x)}{|h|^{d/p+s}}
 =\frac{u(x-h)-u(x)}{|h|^{d/p+s}},
\]
在 \(h=0\) 时补零。差分积分公式说明 \(D_su\in L^p(\mathbb R^{2d})\)，其范数为 \([u]_{s,p}\)。而
\[
 D_s(\tau_zu)(x,h)=D_su(x-z,h).
\]
因此 \([\tau_zu-u]_{s,p}\to0\) 就是 \(L^p(\mathbb R^{2d})\) 中沿向量 \((z,0)\) 的平移连续性。对零阶项同样应用\cref{b1:lem:lp-translation-continuity}，得到第一条结论。

卷积的光滑性来自\secref{b1:sec:2003}。差分与积分相交换，给出
\[
 D_su_\varepsilon(x,h)
 =\int\rho_\varepsilon(z)D_su(x-z,h)\,\mathrm dz
\]
几乎处处成立；其绝对可积性可先用概率测度上的 Hölder 不等式，再用 Tonelli 定理对 \((x,h)\) 积分确认。分别对零阶项及上述差分项使用同一估计，得到
\[
 \|u_\varepsilon-u\|_{W^{s,p}}^p
 \leq\int\rho_\varepsilon(z)\|\tau_zu-u\|_{W^{s,p}}^p\,\mathrm dz.
\]
右侧只使用 \(|z|\leq\varepsilon\) 的平移误差，故由第一条结论趋于零。不减去 \(u\) 时，同一估计与平移不改变范数的事实给出 \(\|u_\varepsilon\|_{W^{s,p}}\leq\|u\|_{W^{s,p}}\)。
\end{proof}

\begin{theorem}\label{b1:thm:fractional-sobolev-smooth-density}
空间 \(C_c^\infty(\mathbb R^d)\) 在 \(W^{s,p}(\mathbb R^d)\) 中稠密。
\end{theorem}
\begin{proof}
先证明可以截去无穷远处的尾部。取 \(0\leq\chi\leq1\)、\(\chi\in C_c^\infty\bigl(B(0,2)\bigr)\)，使 \(\chi=1\) 于 \(\overline B(0,1)\)，并置 \(\chi_R(x)=\chi(x/R)\)、\(R\geq1\)。令 \(w_R=(1-\chi_R)u\)。零阶范数 \(\|w_R\|_p\) 由控制收敛定理趋于零。对差分，写成
\[
 \Delta_hw_R(x)
 =\bigl(1-\chi_R(x)\bigr)\Delta_hu(x)
   +u(x-h)\bigl(\chi_R(x)-\chi_R(x-h)\bigr).
\]
第一项除以 \(|h|^{d/p+s}\) 后，在 \((x,h)\) 上几乎处处趋于零，其 \(p\) 次方受 \(|D_su|^p\) 控制，因此相应积分趋零。

设 \(L\) 为 \(\chi\) 的 Lipschitz 常数。第二项的加权 \(p\) 次积分不超过
\[
 \begin{aligned}
 &\int_{\mathbb R^d}\int_{\mathbb R^d}
  |u(x-h)|^p\frac{\min\{L|h|/R,2\}^p}{|h|^{d+sp}}\,\mathrm dx\,\mathrm dh\\
 &=R^{-sp}\|u\|_p^p
      \int_{\mathbb R^d}\frac{\min\{L|z|,2\}^p}{|z|^{d+sp}}\,\mathrm dz.
 \end{aligned}
\]
等式先在 \(x\) 中使用平移不变性，再令 \(h=Rz\)。最后的核积分在原点附近由 \(|z|^{p-d-sp}\) 控制，在远处由 \(2^p|z|^{-d-sp}\) 控制，故有限。因 \(sp>0\)，这一项也趋于零。合并两项并使用 \(|a+b|^p\leq2^{p-1}(|a|^p+|b|^p)\)，得到 \([w_R]_{s,p}\to0\)。所以 \(\chi_Ru\to u\) 于 \(W^{s,p}\)。

给定误差 \(\eta>0\)，先取 \(R\) 使截断误差小于 \(\eta/2\)，再用前一命题磨光 \(\chi_Ru\)，使磨光误差小于 \(\eta/2\)。卷积的支集包含于 \(\operatorname{supp}\chi_R+\overline B(0,\varepsilon)\)，仍为紧集。所得函数属于 \(C_c^\infty\)，三角不等式完成证明。
\end{proof}

截断证明的第二项不能省去：即使 \(u\) 在远处很小，截断函数的两个取值不同仍会参与点对积分。估计中的 \(R^{-sp}\) 把这种非局部误差控制住，而不是把它当成普通的函数值尾部。

\begin{corollary}\label{b1:cor:compact-fractional-smooth-approximation}
设 \(u\in W^{s,p}(\Omega)\) 在内部紧集 \(K\) 外几乎处处为零。对任意包含 \(K\) 的开集 \(V\subseteq\Omega\)，存在 \(u_j\in C_c^\infty(V)\)，使 \(u_j\to u\) 于 \(W^{s,p}(\Omega)\)。
\end{corollary}
\begin{proof}
由\cref{b1:lem:compact-fractional-zero-extension}先补零得到 \(\widetilde u\in W^{s,p}(\mathbb R^d)\)。前面的磨光命题保证 \(\rho_\varepsilon*\widetilde u\to\widetilde u\) 于全空间范数。取 \(\varepsilon<\operatorname{dist}(K,V^c)\)，其支集便留在 \(V\) 内；限制回 \(\Omega\) 不增加范数误差。若 \(K\) 为空，直接取零函数列。
\end{proof}

\begin{corollary}\label{b1:cor:fractional-domain-smooth-density}
对任意开集 \(\Omega\)，空间 \(C^\infty(\Omega)\cap W^{s,p}(\Omega)\) 在 \(W^{s,p}(\Omega)\) 中稠密。
\end{corollary}
\begin{proof}
取\cref{b1:cor:precompact-smooth-partition}给出的光滑单位分解 \((\psi_i)\) 与局部有限开覆盖 \((V_i)\)，其中 \(\psi_i\in C_c^\infty(V_i)\)、\(V_i\Subset\Omega\)。每个 \(\psi_i\) 是有界 Lipschitz 函数，故由乘子命题，\(\psi_i u\in W^{s,p}(\Omega)\)，且在紧集 \(\operatorname{supp}\psi_i\) 外几乎处处为零。给定 \(\eta>0\)，前一推论允许取 \(v_i\in C_c^\infty(V_i)\)，使
\[
 \|v_i-\psi_i u\|_{W^{s,p}(\Omega)}<\eta2^{-i-1}.
\]
局部有限和 \(v=\sum_i v_i\) 在 \(\Omega\) 内光滑。另一方面，由完备性，误差级数 \(\sum_i(v_i-\psi_i u)\) 在 \(W^{s,p}(\Omega)\) 中收敛到某个 \(e\)，且 \(\|e\|_{W^{s,p}}\leq\eta/2\)。在每个内部相对紧开集上，级数只有有限个非零项，故其和等于 \(v-u\)；范数收敛蕴含局部 \(L^p\) 收敛，因而 \(e=v-u\) 几乎处处。于是 \(v\in W^{s,p}(\Omega)\)，并达到所需误差。
\end{proof}

这里的局部有限光滑和未必具有紧支集。因此最后一个推论不能直接改写成 \(C_c^\infty(\Omega)\) 的稠密性；域内光滑逼近与要求逼近函数在边界附近为零，仍是不同的问题。下一节将把本节的乘子、换元和内部补零估计用于边界图表，使分数阶范数成为迹空间的具体描述。

% !TeX root = ../../Book1.tex
\section{Lipschitz 边界上的迹空间}
\label{b1:sec:2303}

迹属于边界上的 \(L^p\)，只说明它的大小可以控制，还没有说明边界不同位置的函数值如何相互约束。一阶内部导数会给边界留下一个分数阶的差分条件。本节证明：当 \(1<p<\infty\) 时，有界 Lipschitz 区域的一阶迹空间恰为 \(W^{1-1/p,p}(\partial\Omega)\)，而且每个这样的边界数据都有线性、有界地选取的内部延拓。

\ProseParagraphBreak

证明按四步推进：先确认边界分数阶范数不依赖图表与分割函数的选择，再把半空间迹估计增强到 \(W^{1-1/p,p}\)；随后在半空间中令磨光尺度随高度趋零，构造迹的右逆；最后用有限图表与截断把这一构造拼回 Lipschitz 区域。也就是
\[
 \text{图表独立性}
 \longrightarrow \text{半空间分数阶迹}
 \longrightarrow \text{半空间右逆}
 \longrightarrow \text{区域上的有限拼接}.
\]

Di Nezza、Palatucci 与 Valdinoci 的《Hitchhiker's guide to the fractional Sobolev spaces》\cite[Section 2, Proposition 3.8]{DiNezza2012Hitchhikers}介绍了这一联系，并用 Fourier 方法证明二次可积情形。本节改用差分、一个一维积分不等式及尺度卷积，直接处理 \(1<p<\infty\)。
% 实读：DiNezza2012Hitchhikers，作者公开副本 arXiv:1104.4345，
% 第2节末迹空间说明（副本第11页），Proposition 3.8及完整证明（副本第19--21页）。
% 该命题的证明限于p=2；第2节末的一般p陈述没有给出这里所需的完整证明。
% 图表组织另据已实读的第5节；本节Hardy--差分与尺度卷积路线在正文独立展开，
% 不把上述p=2证明写成一般p证明的来源，也不调用后续嵌入或高阶边界理论。

\subsection{局部图表示}
\label{b1:subsec:230301}

先设 \(d\ge2\)，记 \(n=d-1\)，并令
\(\sigma=\mathcal H^n\llcorner\partial\Omega\)。取\secref{b1:sec:2203}的有限内外图表：在刚性坐标中，
\[
 F_j(z,t)=(z,\varphi_j(z)+t),\quad
 D_j=B_j\times(-2h_j,2h_j),\quad
 B_j=(-2\rho_j,2\rho_j)^n,
\]
其中 \(F_j(D_j)\cap\Omega=F_j(D_j\cap H_+)\)。边界参数化为
\(\Gamma_j(z)=F_j(z,0)\)。固定非负函数 \(\chi_j\in C_c^\infty(V_j)\)，其中 \(V_j\) 是对应内图表，使
\[
 \sum_{j=1}^N\chi_j(x)=1\quad(x\in\partial\Omega).
\]
这些函数可由前节构造所用的边界分割项取得；内部项在边界上为零。

\ProseParagraphBreak

\cref{b1:thm:graph-area-formula}给出
\[
 \int_{\Gamma_j(B_j)}q\,\mathrm d\sigma
 =\int_{B_j}q\bigl(\Gamma_j(z)\bigr)
       \sqrt{1+|\nabla\varphi_j(z)|^2}\,\mathrm dz
\]
对非负可测 \(q\) 成立。若 \(\varphi_j\) 的 Lipschitz 常数为 \(L_j\)，面积因子几乎处处介于 \(1\) 与 \(\sqrt{1+L_j^2}\) 之间。因此图上与参数域上的零测集相互对应，\(L^p\) 范数也只相差受控常数。

\begin{definition}\label{b1:def:boundary-fractional-sobolev}
给定 \(0<s<1\)、\(1\le p<\infty\) 及 \(f\in L^p(\partial\Omega,\sigma)\)，将
\[
 f_j(z)=\chi_j\bigl(\Gamma_j(z)\bigr)f\bigl(\Gamma_j(z)\bigr)
 \quad(z\in B_j)
\]
在 \(B_j\) 外补零。若每个 \(f_j\) 都属于 \(W^{s,p}(\mathbb R^n)\)，则称 \(f\) 属于\term[bianjie-fenshujie-sobolev-kongjian]{边界分数阶 Sobolev 空间}[fractional Sobolev space on the boundary]，记为 \(W^{s,p}(\partial\Omega)\)，并规定
\[
 \|f\|_{W^{s,p}(\partial\Omega)}
 \coloneqq\left(\sum_{j=1}^N\|f_j\|_{W^{s,p}(\mathbb R^n)}^p\right)^{1/p}.
\]
\end{definition}
\symidx{\(W^{s,p}(\partial\Omega)\)：用局部图表定义的边界分数阶空间}

分割系数不能从定义中直接删去：将图表边缘处不消失的函数补零，可能引入原函数没有的跳跃。这里每个 \(\chi_j\) 的支集都与外图表的边缘保持正距离。

\begin{proposition}\label{b1:prop:trace-chart-independence}
上述空间不依赖有限图表与分割函数的选择，不同选择给出的范数等价。它在所定范数下完备，而且连续嵌入 \(L^p(\partial\Omega,\sigma)\)。
\end{proposition}
\begin{proof}
由 \(f=\sum_j\chi_j f\)、有限和的不等式及面积因子的上下界，
\[
 \|f\|_{L^p(\sigma)}^p
 \le C\sum_j\|f_j\|_{L^p(\mathbb R^n)}^p
 \le C\|f\|_{W^{s,p}(\partial\Omega)}^p.
\]
这也说明定义中的表达式是范数。

考虑另一套图表 \(\widetilde\Gamma_k\) 与分割函数 \(\widetilde\chi_k\)。在两个边界片相交处，过渡坐标
\[
 \Psi_{jk}=\Gamma_j^{-1}\circ\widetilde\Gamma_k
\]
是两个开子集之间的双 Lipschitz 同胚：图参数化为 Lipschitz 映射，其逆就是相应刚性坐标中的横向投影。待比较的第 \(k\) 个局部函数是有限和
\[
 (\widetilde\chi_k f)\circ\widetilde\Gamma_k
 =\sum_j
   (\widetilde\chi_k\circ\widetilde\Gamma_k)(f_j\circ\Psi_{jk}),
\]
其中每项只需在相应的重叠参数域中解释。

两处分割函数的支集相交部分是重叠图表内的紧集。因而可以在重叠参数域选一个光滑紧支集函数，在该紧集上等于 \(1\)。在右端对应项前乘上这个辅助函数并不改变该项，却使它的支集与重叠域边缘保持正距离。依次使用\cref{b1:prop:fractional-bilipschitz-change}、\cref{b1:prop:fractional-lipschitz-multiplier}及\cref{b1:lem:compact-fractional-zero-extension}，得到该项补零后的 \(W^{s,p}(\mathbb R^n)\) 范数不超过 \(C_{jk}\|f_j\|_{W^{s,p}(\mathbb R^n)}\)。有限求和得到一方向的范数界，交换两套图表得到反向界。

最后，若 \(f^{(m)}\) 在边界范数下为 Cauchy 列，第一步保证它在 \(L^p(\sigma)\) 中收敛到某个 \(f\)。每个局部函数 \(f_j^{(m)}\) 在 \(W^{s,p}(\mathbb R^n)\) 中为 Cauchy 列，由\cref{b1:thm:fractional-sobolev-complete}存在极限。其 \(L^p\) 极限必为 \((\chi_j f)\circ\Gamma_j\) 的零延拓，所以 \(f\) 属于边界空间，且 \(f^{(m)}\) 在定义的范数中收敛到 \(f\)。
\end{proof}

\subsection{半空间迹定理}
\label{b1:subsec:230302}

本节在半空间公式中将 \(\operatorname{Tr}_+\) 简记为 \(\operatorname{Tr}\)。以下令 \(H_+=\mathbb R^n\times(0,\infty)\)，并固定 \(1<p<\infty\) 及 \(s=1-1/p\)。估计边界差分时，把两个边界点抬到与其距离相同的高度。两段竖向差需要下面的\term[hardy-budengshi]{Hardy 不等式}[Hardy inequality]。

\begin{lemma}[Hardy]\label{b1:lem:trace-hardy-inequality}
若 \(a\ge0\) 且 \(a\in L^p(0,\infty)\)，其中 \(1<p<\infty\)，则
\[
 \int_0^\infty r^{-p}\left(\int_0^r a(t)\,\mathrm dt\right)^p\mathrm dr
 \le\left(\frac p{p-1}\right)^p\int_0^\infty a(r)^p\,\mathrm dr.
\]
\end{lemma}
\begin{proof}
先设 \(a\) 有界且支集紧含于 \((0,\infty)\)，令 \(A(r)=\int_0^r a(t)\,\mathrm dt\)。在零点附近 \(A=0\)，在充分大的 \(r\) 处 \(A\) 为常数；由于 \(p>1\)，分部积分的两端边界项都为零。记左端为 \(I\)，则
\[
 I=\frac p{p-1}\int_0^\infty
       \left(\frac{A(r)}r\right)^{p-1}a(r)\,\mathrm dr
 \le\frac p{p-1}I^{(p-1)/p}\|a\|_p.
\]
若 \(I>0\)，约去相应幂次即得结论；\(I=0\) 时也成立。对一般 \(a\)，取
\[
 a_m=\min\{a,m\}\mathbf1_{(1/m,m)},
\]
再由单调收敛定理取极限。
\end{proof}

\begin{theorem}[半空间分数阶迹]\label{b1:thm:half-space-fractional-trace}
设 \(n\ge1\)、\(1<p<\infty\)。半空间的迹算子满足
\[
 \operatorname{Tr}:W^{1,p}(H_+)
 \longrightarrow W^{1-1/p,p}(\mathbb R^n),\quad
 \|\operatorname{Tr}u\|_{W^{1-1/p,p}}
 \le C_{n,p}\|u\|_{W^{1,p}(H_+)}.
\]
\end{theorem}
\begin{proof}
先取整个空间的光滑紧支集函数在闭半空间上的限制，记 \(f(x)=u(x,0)\)。给定 \(h=r\omega\)，其中 \(r>0\)、\(|\omega|=1\)，分解
\[
 \begin{split}
 \Delta_{r\omega}f(x)
 ={}&f(x-r\omega)-u(x-r\omega,r)\\
    &+u(x-r\omega,r)-u(x,r)+u(x,r)-f(x).
 \end{split}
\]
令 \(a(t)=\|\partial_tu(\cdot,t)\|_p\)，并令
\(b(t)=\sum_{i=1}^n\|\partial_i u(\cdot,t)\|_p\)。由微积分基本定理、积分的 Minkowski 不等式及平移不变性，两段竖向差的 \(L^p\) 范数各不超过 \(\int_0^r a(t)\,\mathrm dt\)，而水平差满足
\[
 \|u(\cdot-r\omega,r)-u(\cdot,r)\|_p
 \le r\,b(r).
\]
所以
\[
 \|\Delta_{r\omega}f\|_p^p
 \le C_p\left(\int_0^r a(t)\,\mathrm dt\right)^p
       +C_p r^p b(r)^p,
\]
这里沿用\cref{b1:prop:fractional-difference-formula}的记号
\(\Delta_hf=\tau_hf-f\)，不再把 \(\Delta_h\) 临时改作平移算子。

由于 \(n+sp=n+p-1\)，\cref{b1:prop:fractional-difference-formula}与极坐标变换给出
\[
 [f]_{s,p}^p
 =\int_{\mathbb S^{n-1}}\int_0^\infty
       r^{-p}\|\Delta_{r\omega}f\|_p^p\,\mathrm dr\,
       \mathrm d\mathcal H^{n-1}(\omega).
\]
代入上面的三项估计并用\cref{b1:lem:trace-hardy-inequality}，可得
\[
 [f]_{s,p}^p
 \le C_{n,p}\int_0^\infty\bigl(a(t)^p+b(t)^p\bigr)\,\mathrm dt
 \le C_{n,p}\sum_{i=1}^{n+1}\|\partial_i u\|_{L^p(H_+)}^p.
\]
再加上\cref{b1:thm:half-space-lp-trace}的零阶迹估计，就得到所需全范数界。

对一般 \(u\in W^{1,p}(H_+)\)，先偶反射，再在整个空间中作光滑紧支集逼近，最后限制回半空间。这是\secref{b1:sec:2301}构造 \(L^p\) 迹时使用的逼近。上述估计作用于两项之差，使光滑迹成为 \(W^{s,p}(\mathbb R^n)\) 中的 Cauchy 列；其极限在 \(L^p\) 中又等于既有的 \(\operatorname{Tr}u\)。完备性于是给出结论，并保留同一个范数界。
\end{proof}

指数 \(1-1/p\) 已在积分中出现：边界的径向体积因子将差分核变成 \(r^{-p}\)，恰好与竖向积分的 Hardy 估计匹配。若把 \(p=1\) 直接代入，所用 Hardy 常数发散，而 \(s=0\) 也离开本章分数阶双积分定义的范围；这不是把目标空间写成 \(W^{0,1}\) 就能补上的端点论证。

\subsection{迹的满射性与右逆}
\label{b1:subsec:230303}

仅有迹估计还不能说明每个边界函数都能实现。下面先在半空间中令“磨光尺度等于高度”，同时恢复原数据并控制竖向导数；再把所得右逆逐图表搬回、截断并作有限求和，得到一般有界 Lipschitz 区域上的右逆。

\begin{theorem}\label{b1:thm:half-space-trace-right-inverse}
设 \(n\ge1\)、\(1<p<\infty\)，并令 \(s=1-1/p\)。存在有界线性映射
\[
 R_+:W^{s,p}(\mathbb R^n)\longrightarrow W^{1,p}(H_+),
 \quad \operatorname{Tr}(R_+f)=f.
\]
因此半空间分数阶迹算子为满射。
\end{theorem}
\begin{proof}
固定非负 \(\rho\in C_c^\infty(B(0,1))\)，使 \(\int\rho=1\)，并取 \(\eta\in C_c^\infty(\mathbb R)\)，使 \(\eta=1\) 于 \([-1,1]\)，支集包含于 \((-2,2)\)。对 \(t>0\) 规定
\[
 \rho_t(h)=t^{-n}\rho(h/t),\quad
 v(x,t)=(\rho_t*f)(x),\quad
 (R_+f)(x,t)=\eta(t)v(x,t).
\]
由卷积不等式，\(\|v(\cdot,t)\|_p\le\|f\|_p\)，所以零阶项和竖向求导产生的 \(\eta'(t)v(x,t)\) 都有由 \(\|f\|_p\) 控制的全域 \(L^p\) 范数。

对切向导数，核为 \(t^{-n-1}(\partial_i\rho)(h/t)\)；对高度导数，核为
\[
 \partial_t\rho_t(h)=t^{-n-1}\psi(h/t),\quad
 \psi(z)=-n\rho(z)-z\cdot\nabla\rho(z).
\]
这些核都支撑在 \(|h|<t\) 内，绝对值不超过 \(Ct^{-n-1}\)，且积分为零。最后一点对 \(\psi\) 由紧支集函数的分部积分得到。若 \(K_t\) 是其中任一导数核，因而可以写
\[
 (K_t*f)(x)=\int K_t(h)\Delta_hf(x)\,\mathrm dh.
\]
Hölder 不等式给出
\[
 \begin{split}
 |(K_t*f)(x)|^p
 &\le\left(\int|K_t(h)|\,\mathrm dh\right)^{p-1}
       \int|K_t(h)|\cdot|f(x-h)-f(x)|^p\,\mathrm dh\\
 &\le Ct^{-n-p}\int_{|h|<t}|f(x-h)-f(x)|^p\,\mathrm dh.
 \end{split}
\]
先对 \(x\) 积分，再对 \(0<t<2\) 积分，由 Tonelli 定理，
\[
 \begin{split}
 \int_0^2\|K_t*f\|_p^p\,\mathrm dt
 &\le C\int_{|h|<2}\|\Delta_hf\|_p^p
                  \int_{|h|}^2t^{-n-p}\,\mathrm dt\,\mathrm dh\\
 &\le C_{n,p}\int_{\mathbb R^n}
        \frac{\|\Delta_hf\|_p^p}{|h|^{n+p-1}}\,\mathrm dh
 =C_{n,p}[f]_{s,p}^p.
 \end{split}
\]
函数 \(v\) 在 \(t>0\) 处光滑，因而这些可积的经典导数也是弱导数。结合零阶项及 \(\eta'\) 项的估计，得到
\[
 \|R_+f\|_{W^{1,p}(H_+)}
 \le C_{n,p}\|f\|_{W^{s,p}(\mathbb R^n)}.
\]
所有核与截断都已固定，所以该映射线性。最后，近似恒等核保证 \(v(\cdot,t)\to f\) 于 \(L^p\)，且 \(\eta(t)=1\) 于充分小的正高度。由半空间 \(L^p\) 迹的切片刻画，\(\operatorname{Tr}(R_+f)=f\)。
\end{proof}

这里的 \(R_+\) 是迹算子的右逆：先延入半空间，再取迹，恢复边界数据；反向复合不必恢复任意内部函数，因为零迹函数还可以自由加入。

\begin{theorem}[Lipschitz 边界的分数阶迹]\label{b1:thm:lipschitz-fractional-trace}
设 \(d\ge2\)，\(\Omega\subset\mathbb R^d\) 为有界 Lipschitz 区域，且 \(1<p<\infty\)。既有的迹算子增强为有界满射
\[
 \operatorname{Tr}:W^{1,p}(\Omega)
 \longrightarrow W^{1-1/p,p}(\partial\Omega).
\]
它有一个有界线性右逆 \(R_\Omega\)。此外，边界范数等价于最小延拓范数：
\[
 \|f\|_{W^{1-1/p,p}(\partial\Omega)}
 \asymp
 \inf\bigl\{\,\|u\|_{W^{1,p}(\Omega)}
        \mid \operatorname{Tr}u=f\,\bigr\}.
\]
\end{theorem}
\begin{proof}
先证明有界性。对每个固定边界系数 \(\chi_j\)，将 \(\chi_j u\) 拉回 \(D_j\cap H_+\)，再在侧面和顶面之外补零，得到整个半空间上的函数 \(v_j\)。由边界截断的构造，存在固定紧集 \(L_j\Subset D_j\)，使 \(v_j=0\) 几乎处处成立于 \(H_+\setminus L_j\)；\(L_j\) 可以接触 \(t=0\)，但与图表盒的侧面和顶面有正距离。\secref{b1:sec:2203}局部图表延拓证明中的测试分解因而适用于此处。因此
\[
 \|v_j\|_{W^{1,p}(H_+)}
 \le C_j\|u\|_{W^{1,p}(\Omega)}.
\]
由\cref{b1:thm:lipschitz-lp-trace}的局部端点刻画，它的半空间迹是
\((\chi_j\operatorname{Tr}u)\circ\Gamma_j\) 的零延拓。该相容性已经通过跨边界光滑逼近建立，不要求图函数光滑。半空间分数阶迹估计于是逐项给出
\[
 \|\operatorname{Tr}u\|_{W^{1-1/p,p}(\partial\Omega)}^p
 \le C\|u\|_{W^{1,p}(\Omega)}^p.
\]

现给定 \(f\in W^{s,p}(\partial\Omega)\)，其中 \(s=1-1/p\)。令 \(f_j\) 为定义中的局部零延拓，取 \(w_j=R_+f_j\)。记
\[
 K_j=\operatorname{supp}(\chi_j\circ\Gamma_j)\Subset B_j,
\]
并固定 \(\beta_j\in C_c^\infty(D_j)\)，使
\(\beta_j(z,0)=1\) 于 \(K_j\)。函数
\[
 q_j(z,t)=\beta_j(z,t)w_j(z,t),\quad t>0,
\]
在侧面和顶面附近为零。由 \(w_j(\cdot,t)\to f_j\) 于 \(L^p\)，以及 \(\beta_j(\cdot,t)\) 一致趋于 \(\beta_j(\cdot,0)\)，可得
\(\operatorname{Tr}q_j=\beta_j(\cdot,0)f_j=f_j\)。
把 \(q_j\) 经 \(F_j^{-1}\) 推回 \(\Omega\cap F_j(D_j)\)，再在 \(\Omega\) 的其余部分补零，记结果为 \(Q_jf\)。由于 \(\beta_j\in C_c^\infty(D_j)\)，\(q_j\) 在固定紧集 \(\operatorname{supp}\beta_j\) 外几乎处处为零；这一紧集与侧面和顶面有正距离，故局部测试恒等式可以拼接。这里始终只在 \(\Omega\) 内补零，并没有跨越 \(\partial\Omega\) 作零延拓。故
\[
 \|Q_jf\|_{W^{1,p}(\Omega)}
 \le C_j\|f_j\|_{W^{s,p}(\mathbb R^n)}.
\]
迹与局部坐标的相容性给出
\(\operatorname{Tr}(Q_jf)=\chi_j f\)，包括图表外的零部分。

规定
\[
 R_\Omega f=\sum_{j=1}^N Q_jf.
\]
所有图表、边界系数、卷积核与 \(\beta_j\) 都固定不变，所以它是线性的。有限求和与上述估计给出有界性，而
\[
 \operatorname{Tr}(R_\Omega f)
 =\sum_{j=1}^N\chi_j f=f
\]
证明右逆性质和满射性。最后，对任意迹为 \(f\) 的 \(u\)，迹估计给出边界范数不超过 \(C\|u\|_{W^{1,p}}\)；反向取 \(u=R_\Omega f\) 即得最小延拓范数的等价关系。
\end{proof}

这个迹空间允许的边界不连续性取决于 \(p\)。例如在二维半空间上取边界函数 \(f=\mathbf1_{(0,1)}\)。\cref{b1:ex:23-interval-characteristic}的差分计算表明，它属于 \(W^{s,p}(\mathbb R)\) 当且仅当 \(sp<1\)。代入 \(s=1-1/p\)，就知道它能作为 \(W^{1,p}(H_+)\) 函数的迹，当且仅当 \(1<p<2\)。因此不能从内部具有一阶弱导数，直接推断边界数据连续。

\ProseParagraphBreak

本节没有把两个端点纳入分数阶结论。对 \(p=1\)，\secref{b1:sec:2301}仍给出取值于 \(L^1(\partial\Omega)\) 的连续迹与零迹刻画；本章不据此声称 \(L^1(\partial\Omega)\) 是迹的精确像空间，也不声称该迹满射或具有同样的线性右逆。对 \(p=\infty\)，既不能代入有限 \(p\) 的双积分范数，也不能用该范数的稠密性取极限；应回到局部 Lipschitz 代表与边界几何讨论。

\ProseParagraphBreak

若 \(d=1\)，有界 Lipschitz 开集的每个边界点都在边界中孤立，边界又紧，故只有有限多个端点。迹空间此时就是这些端点上的有限维数据空间：在每个区间分量上作连接两端数据的仿射函数，便得到线性右逆，其范数由有限多个区间长度控制。这个论证也处理相应端点指数，但不把 \(n=0\) 代入上面的欧氏差分核。

% !TeX root = ../../Book1.tex
\OptionalSection{高阶迹理论}{b1:sec:2304}

一阶迹决定函数在边界上的零阶数据。若域内还有更多弱导数，便可以继续对这些导数取迹；但所得数据不是一组互不相关的边界函数。切向导数必须与原函数迹的切向变化一致，法向导数则承担另外的边界信息。本节先在平坦边界上证明这种结构，再说明弯曲边界的几何限制。完整的高阶边界数据延拓不在本节证明范围内。

\subsection{\texorpdfstring{\(W^{k,p}\)}{Wk,p} 的边界数据}
\label{b1:subsec:230401}

设 \(d\ge2\)、\(n=d-1\)、\(H_+=\mathbb R^n\times(0,\infty)\)，并取整数 \(k\ge1\)、\(1<p<\infty\)。若 \(u\in W^{k,p}(H_+)\)，那么对每个 \(|\alpha|\le k-1\)，都有 \(\partial^\alpha u\in W^{1,p}(H_+)\)，所以
\[
 g_\alpha=\operatorname{Tr}_+(\partial^\alpha u)
\]
已经有明确定义。各个 \(g_\alpha\) 至少属于 \(W^{s,p}(\mathbb R^n)\)，其中 \(s=1-1/p\)；这是把一阶分数迹定理用于 \(\partial^\alpha u\) 的直接结果。下面还要利用不同阶数据之间的相容关系。

先约定非整数高阶空间的记号。对整数 \(m\ge0\)、\(0<s<1\)，令
\[
 W^{m+s,p}(\mathbb R^n)
 =\{\,f\in W^{m,p}(\mathbb R^n)
       \mid \partial^\beta f\in W^{s,p}(\mathbb R^n)
               \text{ 对所有 }|\beta|=m\,\},
\]
其范数取为
\[
 \|f\|_{W^{m+s,p}}^p
 =\|f\|_{W^{m,p}}^p+
     \sum_{|\beta|=m}[\partial^\beta f]_{s,p}^p.
\]
当 \(m=0\) 时，这正是\secref{b1:sec:2302}的定义。这里对最高整数阶弱导数增加分数阶差分条件，而不是把原先的整数阶定义中的“阶数”形式地换成实数。

\begin{proposition}[平坦边界的高阶迹]\label{b1:prop:flat-higher-traces}
设 \(u\in W^{k,p}(H_+)\)、\(1<p<\infty\)。对 \(|\alpha|\le k-2\) 及切向方向 \(1\le i\le n\)，有
\begin{equation}\label{b1:eq:trace-tangential-compatibility}
 \partial_i g_\alpha=g_{\alpha+e_i}
 \quad\text{作为 }\mathbb R^n\text{ 上的弱导数}.
\end{equation}
若记
\[
 \gamma_j u=\operatorname{Tr}_+(\partial_d^j u),
 \quad 0\le j\le k-1,
\]
则
\begin{equation}\label{b1:eq:higher-flat-trace-bound}
 \gamma_j u\in W^{k-j-1/p,p}(\mathbb R^n),\qquad
 \|\gamma_j u\|_{W^{k-j-1/p,p}}
 \le C_{d,k,p}\|u\|_{W^{k,p}(H_+)}.
\end{equation}
\end{proposition}
\begin{proof}
先证切向相容性。固定 \(|\alpha|\le k-2\)，令 \(v=\partial^\alpha u\)、\(w=\partial_i v\)。两者都属于 \(W^{1,p}(H_+)\)。由 Fubini 定理及切向弱导数恒等式，对几乎每个 \(t>0\)，切片满足
\[
 \int_{\mathbb R^n}v(x',t)\partial_i\phi(x')\,\mathrm dx'
 =-\int_{\mathbb R^n}w(x',t)\phi(x')\,\mathrm dx'
 \quad\bigl(\phi\in C_c^\infty(\mathbb R^n)\bigr).
\]
这里可以先对可数稠密的测试函数族取得共同的好切片，再由局部积分连续性扩展到全部测试函数。选取这样的 \(t_\ell\downarrow0\)。由\eqref{b1:eq:trace-slice-convergence}，\(v(\cdot,t_\ell)\to g_\alpha\)、\(w(\cdot,t_\ell)\to g_{\alpha+e_i}\) 于 \(L^p(\mathbb R^n)\)。在积分恒等式中取极限，便得\eqref{b1:eq:trace-tangential-compatibility}。

令 \(m=k-1-j\)。反复使用切向相容性，得到
\[
 \partial_{x'}^\beta\gamma_j u
 =\operatorname{Tr}_+
       (\partial_{x'}^\beta\partial_d^j u)
 \quad(|\beta|\le m).
\]
右侧由半空间的 \(L^p\) 迹估计控制，因为内部导数总阶数至多 \(k-1\)，还能再取一阶弱导数。这证明 \(\gamma_ju\in W^{m,p}(\mathbb R^n)\) 并控制其整数阶范数。

对 \(|\beta|=m\)，内部函数
\(\partial_{x'}^\beta\partial_d^j u\in W^{1,p}(H_+)\)。再用\cref{b1:thm:half-space-fractional-trace}，得到它的迹属于 \(W^{s,p}\)，且相应半范数受 \(\|u\|_{W^{k,p}}\) 控制。有限多个 \(\beta\) 求和，即得 \(\gamma_ju\in W^{m+s,p}\) 及所需估计，因为 \(m+s=k-j-1/p\)。
\end{proof}

例如 \(u\in W^{2,p}(H_+)\) 时，零阶迹属于 \(W^{2-1/p,p}\)，法向方向的一阶导数迹属于 \(W^{1-1/p,p}\)，而每个切向导数迹必须等于零阶迹的对应弱导数。不能先任意指定函数值及全部一阶偏导数的边界值，再期待存在一个内部函数实现它们。

上述证明只需对已有整数阶导数重复使用一阶迹，不依赖高阶函数在闭半空间上的光滑稠密性。它给出的是连续性和数据相容性；同时为所有 \(\gamma_j\) 构造有界右逆，还需要能匹配多项边界数据的延拓公式，不能把零阶的尺度卷积原样使用 \(k\) 次就视为完成。

\subsection{法向导数}
\label{b1:subsec:230402}

在上半空间的边界上，外单位法向为 \(\nu=-e_d\)。因此对 \(u\in W^{2,p}(H_+)\)，\term[faxiang-daoshu]{法向导数}[normal derivative]的边界迹为
\[
 \partial_\nu u=-\operatorname{Tr}_+(\partial_du)=-\gamma_1u.
\]
更高阶沿固定法向的导数相应带有 \((-1)^j\)。本节的 \(\gamma_j\) 始终按正 \(x_d\) 方向定义；写外法向数据时另加这个符号，避免把方向约定隐藏在记号中。

设 \(\Omega\) 为有界 Lipschitz 区域。边界图表几乎处处给出外单位法向 \(\nu\)，其分量本性有界。对 \(u\in W^{2,p}(\Omega)\)，各一阶导数属于 \(W^{1,p}(\Omega)\)，所以可以在 \(L^p(\partial\Omega)\) 中定义
\[
 \partial_\nu u
 =\sum_{i=1}^d\nu_i\operatorname{Tr}(\partial_i u).
\]
由 \(|\nu|=1\) 及 \(L^p\) 迹估计，这是受 \(\|u\|_{W^{2,p}(\Omega)}\) 控制的边界函数。若 \(u\) 与边界足够光滑，它与通常的方向导数相同。

但是，本性有界乘子未必保持分数阶 Sobolev 空间。\secref{b1:sec:2302}的乘子结论还要求 Lipschitz 控制；在折角处，\(\nu\) 可以跳跃。因此不能仅凭每个 \(\operatorname{Tr}(\partial_i u)\) 属于 \(W^{1-1/p,p}(\partial\Omega)\)，就断言它们与 \(\nu_i\) 相乘后仍属于该空间。

一个具体例子是正方形 \(\Omega=(0,1)^2\) 与光滑函数 \(u(x,y)=x\)。在下边上 \(\partial_\nu u=0\)，在左边上 \(\partial_\nu u=-1\)。在原点附近用折线弧长作参数，这个边界函数有一个大小为 \(1\) 的跳跃。若它属于 \(W^{s,p}\)，其半范数必须包含有限的跨角点积分；然而该积分从下方控制
\[
 \int_0^\varepsilon\int_0^\varepsilon
       \frac{\mathrm da\,\mathrm db}{(a+b)^{1+sp}}.
\]
在 \(0<a<\varepsilon/2\)、\(a<b<2a\) 上积分，便得到常数倍的
\(\int_0^{\varepsilon/2}a^{-sp}\,\mathrm da\)，当 \(sp\ge1\) 时发散。因此取 \(s=1-1/p\)、\(p\ge2\)，这个法向数据不属于预期的分数阶空间，尽管 \(u\) 在整个平面上光滑。

若边界具有更多光滑性，法向及坐标过渡也具有更多可控导数，平坦边界的高阶估计才有望在图表之间保持。对高阶法向导数，还须说明是在边界点固定方向作高阶导数，还是把法向场延到邻域后反复作用；后者可能引入法向场的导数。高阶迹理论同时涉及函数正则性、边界正则性和数据约定，三者不能只靠一个“边界光滑”的笼统表述代替。

\subsection{Besov 空间作为自然迹空间的预告}
\label{b1:subsec:230403}

一阶迹的指数 \(1-1/p\)，来自对法向厚度的积分与切向差分尺度之间的配合。允许更一般的光滑阶数和尺度求和方式后，便进入\term[besov-kongjian]{Besov 空间}[Besov space]。这里只介绍本章已经能够理解的一种差分定义，不预用 Fourier 变换。

先设 \(0<s<1\)、\(1\le p<\infty\) 及 \(1\le q<\infty\)。对 \(f\in L^p(\mathbb R^n)\)，令
\[
 [f]_{B^s_{p,q}}
 =\left(
   \int_{\mathbb R^n}
    \frac{\|\Delta_hf\|_{L^p}^q}{|h|^{sq}}
    \,\frac{\mathrm dh}{|h|^n}
   \right)^{1/q}.
\]
要求这个量有限，并以 \(\|f\|_p+[f]_{B^s_{p,q}}\) 为范数，得到本节所用的非齐次 Besov 空间 \(B^s_{p,q}(\mathbb R^n)\)。另外，当 \(q=\infty\) 时规定
\[
 [f]_{B^s_{p,\infty}}
 =\operatorname*{ess\,sup}_{h\ne0}
       \frac{\|\Delta_hf\|_{L^p}}{|h|^s}.
\]
由于大尺度差分受 \(2\|f\|_p\) 控制，也可以只在 \(|h|<1\) 上积分或取本性上确界，得到等价范数。

两个指数承担不同的平均：\(p\) 先度量固定平移 \(h\) 下、空间变量中的差分大小；\(q\) 再度量这些大小随平移尺度的分布。当 \(q=p\) 时，Tonelli 定理直接给出
\[
 [f]_{B^s_{p,p}}^p
 =\int_{\mathbb R^n}\int_{\mathbb R^n}
       \frac{|\Delta_hf(x)|^p}{|h|^{n+sp}}\,\mathrm dx\,\mathrm dh
 =[f]_{s,p}^p.
\]
因此在当前 \(0<s<1\) 的范围内，
\[
 B^s_{p,p}(\mathbb R^n)=W^{s,p}(\mathbb R^n)
\]
作为集合相同，所列非齐次范数等价。一阶迹空间也就可以写成 \(B^{1-1/p}_{p,p}\)，并不是又多出了一种不同的边界条件。

对更高的非整数阶数，需要使用高阶差分或先取整数阶导数；这些定义之间的等价、一般 Besov 空间的插值以及完整高阶迹右逆，属于进一步的函数空间理论。可从 Triebel 的《Theory of Function Spaces》\cite{Triebel1983Theory}继续阅读，也可在 Adams 与 Fournier 的《Sobolev Spaces》\cite{Adams2003Sobolev}中衔接经典 Sobolev 理论。这里不将尚未展开的一般 Besov 迹定理用作后续正文的前置。

至此，边界值的定义、其一阶分数正则性及零边界闭包已连成一条证明链。下一章回到区域内部：在这些逼近与延拓工具准备好以后，将研究梯度怎样控制函数的振幅和可积性。这些嵌入估计与本章的迹估计相互补充，但处理的是不同的测度和不同位置上的函数信息。


\begin{chaptersummary}
迹由内部沿线绝对连续性和 Sobolev 范数估计确定，而不是任意边界赋值。对有界 Lipschitz 区域，有限指数下迹的核正是 \(W_0^{1,p}\)；反向证明必须在图表中控制支集，再将不一定光滑的坐标推回函数作内部逼近。

分数阶范数比较所有点对，因而能检测跨越边界或连通分支的差异。乘子、补零、截断和坐标变换都要控制这种非局部比较。对于 \(1<p<\infty\)，法向积分和切向尺度配合给出 \(1-1/p\) 阶边界正则性；高度随磨光尺度变化的构造，又把任意这样的数据延回内部。

高阶数据进一步具有切向相容性。平坦边界的结论不能不加条件地转移到有折角的法向场上。本章因此把已证明的一阶边界理论、平坦高阶估计和进一步的 Besov 理论分开。接下来的嵌入与紧性，将主要利用前两章和本章已经完成的逼近、延拓及零边界工具。
\end{chaptersummary}

\BookExercises

% \chapref{b1:ch:23}习题临时稿；不另设 BookExercises 入口。
% 八题均为自拟的工具变式；来源为已建立的正文接口与此处独立展开的计算，不冒认教材原题编号。
% 仅用本章及前章工具，不使用后续嵌入、Fourier或尚未建立的Besov表示定理。

\begin{exercise}\label{b1:ex:23-interval-characteristic}
设 \(0<s<1\)、\(1\le p<\infty\)，令 \(u=\mathbf1_{(0,1)}\)，把它看作整个 \(\mathbb R\) 上的函数。证明
\[
 u\in W^{s,p}(\mathbb R)\quad\Longleftrightarrow\quad sp<1,
\]
并在有限情形计算半范数 \([u]_{s,p}^p\) 的精确值。说明临界条件 \(sp=1\) 不能归入有限的一侧。
\end{exercise}
% 来源：自拟；直接计算定义中的跨跳跃双积分，区别于原区间内常函数的零半范数。
\begin{solution}
函数的 \(L^p\) 范数为 \(1\)，问题完全在于半范数。只有一个变量在 \((0,1)\)、另一个变量在其外时，函数差才非零。由对称性及 Tonelli 定理，
\[
 \begin{aligned}
 [u]_{s,p}^p
 &=2\int_0^1\left(
     \int_{-\infty}^0\frac{\mathrm dy}{(x-y)^{1+sp}}
    +\int_1^\infty\frac{\mathrm dy}{(y-x)^{1+sp}}
             \right)\mathrm dx\\
 &=\frac2{sp}\int_0^1
       \bigl(x^{-sp}+(1-x)^{-sp}\bigr)\,\mathrm dx.
 \end{aligned}
\]
当 \(sp<1\) 时，两个端点积分各等于 \(1/(1-sp)\)，故
\[
 [u]_{s,p}^p=\frac4{sp(1-sp)}.
\]
若 \(sp\ge1\)，至少一个端点积分发散；在 \(sp=1\) 时就是 \(\int_0^1x^{-1}\,\mathrm dx\) 的对数发散。因此条件为严格不等式。

若仅在 \((0,1)\) 内计算，\(u\) 是常函数，双积分为零。全空间半范数还比较边界两侧的函数值，正是这些跨越端点的配对检测到了跳跃。
\end{solution}

\begin{exercise}\label{b1:ex:23-smooth-integer-limit}
设 \(d\ge2\)、\(1\le p<\infty\)，实值 \(u\in C_c^\infty(\mathbb R^d)\)。记 \(\mathrm d\sigma\) 为单位球面上的 \(\mathcal H^{d-1}\) 测度。证明
\[
 \lim_{s\uparrow1}(1-s)[u]_{s,p}^p
 =\frac1p\int_{\mathbb S^{d-1}}\int_{\mathbb R^d}
            |\nabla u(x)\cdot\omega|^p\,\mathrm dx\,\mathrm d\sigma(\omega).
\]
特别地，\(p=2\) 时右端等于
\[
 \frac{\sigma(\mathbb S^{d-1})}{2d}
       \int_{\mathbb R^d}|\nabla u|^2\,\mathrm dx.
\]
本题只要求证明光滑紧支集情形，不把所得极限未经说明地推广到一般函数。
\end{exercise}
% 来源：自拟；差分公式、Taylor余项与径向积分的组合计算，不预借一般整数阶极限定理。
\begin{solution}
由\cref{b1:prop:fractional-difference-formula}，半范数可写成平移差的积分。先估计 \(|h|\ge1\) 的部分。平移等距性与三角不等式给出
\[
 \|u(\,\cdot+h)-u\|_p^p\le2^p\|u\|_p^p,
\]
所以
\[
 (1-s)\int_{|h|\ge1}
 \frac{\|u(\,\cdot+h)-u\|_p^p}{|h|^{d+sp}}\,\mathrm dh
 \le\frac{(1-s)2^p\sigma(\mathbb S^{d-1})}{sp}\|u\|_p^p
 \longrightarrow0.
\]

对 \(0<r<1\)、\(\omega\in\mathbb S^{d-1}\)，记
\[
 A(r,\omega)=
 \left\|\frac{u(\,\cdot+r\omega)-u}{r}\right\|_p^p,
 \quad A(0,\omega)=\|\nabla u\cdot\omega\|_p^p.
\]
微积分基本定理给出
\[
 \frac{u(x+r\omega)-u(x)}r
 =\int_0^1\nabla u(x+tr\omega)\cdot\omega\,\mathrm dt.
\]
因为 \(\nabla u\) 一致连续且有紧支集，右端与
\(\nabla u(x)\cdot\omega\) 之差在 \(L^p(\mathbb R^d)\) 中趋零，而且对全部单位向量 \(\omega\) 一致：函数差的支集包含在一个固定有界集内，其上确界受 \(\nabla u\) 在距离不超过 \(r\) 的两点间的振幅控制。各差商的 \(L^p\) 范数又一致有界，故
\[
 \sup_\omega|A(r,\omega)-A(0,\omega)|\longrightarrow0.
\]

令 \(B(r)=\int_{\mathbb S^{d-1}}A(r,\omega)\,\mathrm d\sigma(\omega)\)，则 \(B\) 有界且 \(B(r)\to B(0)\)。用径向积分处理 \(|h|<1\) 的部分，得到
\[
 (1-s)\int_0^1r^{p(1-s)-1}B(r)\,\mathrm dr.
\]
这段权重的总积分为 \(1/p\)。给定 \(0<\delta<1\)，在 \((0,\delta)\) 上用 \(B(r)\) 接近 \(B(0)\) 控制误差；在 \((\delta,1)\) 上，权重质量为
\[
 \frac{1-\delta^{p(1-s)}}p\longrightarrow0.
\]
先选小 \(\delta\)，再令 \(s\uparrow1\)，便知所列积分趋于 \(B(0)/p\)。结合远处部分趋零，得到第一条公式。

当 \(p=2\) 时，球面对称性给出
\(\int\omega_i\omega_j\,\mathrm d\sigma=0\)（\(i\ne j\)），各
\(\int\omega_i^2\,\mathrm d\sigma\) 相等，而它们之和为
\(\sigma(\mathbb S^{d-1})\)。展开 \(|\nabla u\cdot\omega|^2\) 后逐项积分，便得第二条公式。有限差分在小尺度下读出方向导数，系数 \(1/p\) 来自径向权重，而不是可以任意省略的归一化。
\end{solution}

\begin{exercise}\label{b1:ex:23-fractional-modulus-truncation}
设 \(0<s<1\)、\(1\le p<\infty\)，实值 \(u\in W^{s,p}(\Omega)\)。
\begin{enumerate}
\item\label{b1:item:23-fractional-modulus}
证明 \(|u|\in W^{s,p}(\Omega)\)，且 \([\,|u|\,]_{s,p;\Omega}\le[u]_{s,p;\Omega}\)。当 \(p=2\) 时，求两个半范数平方之差，并证明半范数相等当且仅当 \(u\) 几乎处处不变号。
\item\label{b1:item:23-fractional-truncation}
令 \(T_M(t)=\max\bigl(-M,\min(t,M)\bigr)\)。证明 \(T_M(u)\to u\) 于 \(W^{s,p}(\Omega)\)。
\end{enumerate}
\end{exercise}
% 来源：自拟；非线性收缩、变号能量损失与可积双积分上的控制收敛，不重复正文线性乘子及光滑稠密性。
\begin{solution}
由反三角不等式，
\[
 \bigl||u(x)|-|u(y)|\bigr|\le|u(x)-u(y)|.
\]
将其 \(p\) 次幂除以核的分母并积分，即得半范数估计；零阶 \(L^p\) 范数则完全相等。

若 \(p=2\)，对任意实数 \(a,b\)，
\[
 |a-b|^2-\bigl||a|-|b|\bigr|^2
 =4|a|\cdot|b|\,\mathbf1_{\{ab<0\}}.
\]
两边积分得到
\[
 [u]_{s,2;\Omega}^2-[\,|u|\,]_{s,2;\Omega}^2
 =4\iint_{\{u(x)u(y)<0\}}
       \frac{|u(x)|\cdot|u(y)|}{|x-y|^{d+2s}}\,\mathrm dx\,\mathrm dy.
\]
两个半范数都有限，所以这里没有无穷量相减。若 \(u\) 几乎处处不变号，右侧为零。反之，若正值集和负值集都具有正测度，则其乘积集具有正的乘积测度，核上的被积函数在那里严格为正，故积分严格为正。这证明等号的刻画。

对于第二问，记 \(R_M(t)=t-T_M(t)\)。这个实函数连续、分段线性，各段斜率为 \(0\) 或 \(1\)，因此
\[
 |R_M(a)-R_M(b)|\le|a-b|,\quad |R_M(a)|\le|a|.
\]
又因为有限实数 \(a\) 满足 \(R_M(a)\to0\)，控制收敛先给出
\(\|R_M(u)\|_p\to0\)。对于双积分，几乎每一对 \((x,y)\) 的分子也趋于零，并且
\[
 \frac{\bigl|R_M\bigl(u(x)\bigr)-R_M\bigl(u(y)\bigr)\bigr|^p}{|x-y|^{d+sp}}
 \le\frac{|u(x)-u(y)|^p}{|x-y|^{d+sp}}.
\]
右端按假设可积，再次应用控制收敛，得到
\([R_M(u)]_{s,p;\Omega}\to0\)。这就是所求范数收敛。两处收敛都使用有限 \(p\) 的积分控制，并非对所有函数值统一取上确界。
\end{solution}

\begin{exercise}\label{b1:ex:23-boundary-concentration}
设 \(d\ge2\)、\(1\le p<\infty\)，并记
\(\Omega=(0,1)^{d-1}\times(0,1)\)、\(B'=(0,1)^{d-1}\)。
取非零 \(f\in C_c^\infty(B')\) 及满足 \(\chi(0)=1\)、\(\operatorname{supp}\chi\subset(-1,1)\) 的 \(\chi\in C_c^\infty(\mathbb R)\)。对 \(0<\varepsilon<1\)，令
\[
 u_\varepsilon(x',t)=f(x')\chi(t/\varepsilon).
\]
计算 \(u_\varepsilon\) 的体积 \(L^p\) 范数、边界迹的 \(L^p\) 范数以及
\(\|\partial_tu_\varepsilon\|_{L^p(\Omega)}\) 随 \(\varepsilon\) 的变化。由此证明：不存在与 \(u\) 无关的常数 \(C\)，使所有在 \(\overline\Omega\) 附近光滑的函数都满足
\[
 \|\operatorname{Tr}u\|_{L^p(\partial\Omega)}
 \le C\|u\|_{L^p(\Omega)}.
\]
说明这个例子为什么不反驳本章的 \(W^{1,p}\) 迹估计。
\end{exercise}
% 来源：自拟；边界层缩放例，用来分开体积收敛、迹收敛与一阶导数控制。
\begin{solution}
记两个严格为正的常数
\[
 A_p=\left(\int_0^1|\chi(r)|^p\,\mathrm dr\right)^{1/p},
 \quad
 B_p=\left(\int_0^1|\chi'(r)|^p\,\mathrm dr\right)^{1/p}.
\]
\(A_p>0\) 来自 \(\chi(0)=1\) 及连续性；\(B_p>0\) 是因为 \(\chi\) 在 \(1\) 的邻域为零，不可能在 \([0,1]\) 上恒为常数。代换 \(t=\varepsilon r\) 得
\[
 \|u_\varepsilon\|_{L^p(\Omega)}
   =\varepsilon^{1/p}A_p\|f\|_{L^p(B')},
 \quad
 \|\partial_tu_\varepsilon\|_{L^p(\Omega)}
   =\varepsilon^{1/p-1}B_p\|f\|_{L^p(B')}.
\]
底面上的限制为 \(f\)，顶面上的限制为零；由于 \(f\) 的支集包含在 \(B'\) 的内部，各侧面上的限制也为零。棱边的 \(\mathcal H^{d-1}\) 测度为零，故光滑函数的迹满足
\[
 \|\operatorname{Tr}u_\varepsilon\|_{L^p(\partial\Omega)}
 =\|f\|_{L^p(B')}>0.
\]
若题设估计成立，约去这个正数就会得到
\(1\le C A_p\varepsilon^{1/p}\)，令 \(\varepsilon\downarrow0\) 即矛盾。

另一方面，法向导数的范数在 \(p>1\) 时趋于无穷，在 \(p=1\) 时保持为严格正的常数；切向导数则满足
\[
 \|\partial_{x_i}u_\varepsilon\|_{L^p(\Omega)}
 =\varepsilon^{1/p}A_p\|\partial_{x_i}f\|_{L^p(B')}.
\]
因此 \(u_\varepsilon\) 虽然在体积 \(L^p\) 中趋于零，却不在 \(W^{1,p}\) 中趋于零。迹估计使用的一阶导数项正好检测到了这层越来越薄的过渡区域。
\end{solution}

\begin{exercise}\label{b1:ex:23-trace-right-inverse-freedom}
设 \(\Omega\subset\mathbb R^d\) 为非空有界 Lipschitz 开集，\(d\ge2\)、\(1<p<\infty\)。记
\[
 X=W^{1,p}(\Omega),\quad
 Y=W^{1-1/p,p}(\partial\Omega),\quad
 T=\operatorname{Tr}:X\longrightarrow Y,
\]
并取\cref{b1:thm:lipschitz-fractional-trace}给出的一个有界线性右逆 \(R:Y\to X\)。
\begin{enumerate}
 \item\label{b1:item:23-trace-right-inverse-difference}
 证明有界线性算子 \(S:Y\to X\) 也是 \(T\) 的右逆，当且仅当
 \(S=R+K\)，其中 \(K:Y\to W_0^{1,p}(\Omega)\) 是有界线性算子。实际构造一个非零的 \(K\)，说明右逆一般不唯一。
 \item\label{b1:item:23-trace-affine-class}
 对任意 \(g\in Y\)，证明
 \[
  \{u\in X:Tu=g\}=Rg+W_0^{1,p}(\Omega).
 \]
 说明为什么用另一个右逆替换 \(R\)，不会改变给定边界值的这一函数类。
\end{enumerate}
\end{exercise}
% 来源：自拟；把迹满射和零迹刻画组合为仿射 Dirichlet 类，不预用后续偏微分方程存在定理。
\begin{solution}
由\cref{b1:thm:zero-trace-characterization}，
\(\ker T=W_0^{1,p}(\Omega)\)。若 \(TS=\operatorname{id}_Y=TR\)，则
\[
 T(S-R)=0,
\]
所以 \(K=S-R\) 的像包含于 \(W_0^{1,p}(\Omega)\)。这个子空间采用 \(X\) 的限制范数，因此 \(K\) 作为取值于该子空间的算子仍然有界。反过来，若 \(K\) 的像包含于 \(\ker T\)，则 \(T(R+K)=\operatorname{id}_Y\)。

下面给出非零的选择。记 \(\sigma=\mathcal H^{d-1}|_{\partial\Omega}\)，取非零 \(w\in C_c^\infty(\Omega)\)，并令
\[
 \ell(g)=\int_{\partial\Omega}g\,\mathrm d\sigma,
 \quad K(g)=\ell(g)w.
\]
有界 Lipschitz 开集的边界具有有限且严格为正的 \(\sigma\) 测度，而本章选定的边界范数控制 \(L^p(\partial\Omega)\) 范数。因此 Hölder 不等式给出
\[
 |\ell(g)|\le
 \sigma(\partial\Omega)^{1-1/p}\|g\|_{L^p(\partial\Omega)}
 \le C\|g\|_Y.
\]
常函数 \(1\) 是 \(X\) 中常函数的迹，故属于 \(Y\)，并且
\(\ell(1)=\sigma(\partial\Omega)>0\)。于是 \(K\ne0\)，且其像包含于
\(\operatorname{span}\{w\}\subset W_0^{1,p}(\Omega)\)。这便构造出两个不同的右逆 \(R\) 与 \(R+K\)。

最后，\(Tu=g\) 等价于 \(T(u-Rg)=0\)，也就等价于
\(u-Rg\in W_0^{1,p}(\Omega)\)，这证明第二问的集合恒等式。如果换用 \(S\)，则
\(Sg-Rg\in W_0^{1,p}(\Omega)\)；用子空间中的一个元素平移这个子空间不会改变所得陪集。因此边界条件本身是确定的，选择右逆只是选择了其中一个起点。
\end{solution}

\begin{exercise}\label{b1:ex:23-circle-energy}
把单位圆盘 \(D\subset\mathbb R^2\) 记为复平面中的 \(\{z:|z|<1\}\)。对 \(n\ge1\)，令
\[
 U_n(z)=z^n,\quad g_n(e^{i\theta})=e^{in\theta}.
\]
本题单独采用圆周上的弦长双积分
\[
 Q(g)=\int_{\mathbb S^1}\int_{\mathbb S^1}
        \frac{|g(x)-g(y)|^2}{|x-y|^2}\,
                 \mathrm d\mathcal H^1(x)\,\mathrm d\mathcal H^1(y).
\]
证明
\[
 \|g_n\|_{L^2(\mathbb S^1)}^2=2\pi,\quad
 Q(g_n)=4\pi^2 n,\quad
 \int_D|\nabla U_n|^2\,\mathrm dx=2\pi n,
\]
其中复值函数的梯度平方是
\(|\partial_xU_n|^2+|\partial_yU_n|^2\)。
只用有限几何级数与积分计算完成证明，不使用 Fourier 展开定理。
\end{exercise}
% 来源：自拟；圆周高频数据和显式光滑延拓的精确计算。Q 为题内指定双积分，不冒认为正文固定图表范数的原值。
\begin{solution}
因为 \(|g_n|=1\)，第一个等式直接由圆周长度得到。在几乎每一对参数 \((\theta,\varphi)\) 处，有限几何级数给出
\[
 \frac{|e^{in\theta}-e^{in\varphi}|^2}
      {|e^{i\theta}-e^{i\varphi}|^2}
 =\left|\sum_{j=0}^{n-1}e^{ij(\theta-\varphi)}\right|^2.
\]
对于整数 \(j,\ell\)，初等积分满足
\[
 \int_0^{2\pi}e^{i(j-\ell)h}\,\mathrm dh
 =
 \begin{cases}
  2\pi,&j=\ell,\\
  0,&j\ne\ell.
 \end{cases}
\]
展开有限和的模平方，再逐项积分，得到
\[
 \int_0^{2\pi}\left|\sum_{j=0}^{n-1}e^{ijh}\right|^2
                        \,\mathrm dh=2\pi n.
\]
对固定的 \(\varphi\)，以 \(h=\theta-\varphi\) 代换并用周期性，内层积分仍等于 \(2\pi n\)。外层再积分一次便得 \(Q(g_n)=4\pi^2n\)。

在圆盘内直接求导，
\[
 \partial_x U_n=n z^{n-1},\quad
 \partial_y U_n=in z^{n-1}.
\]
所以两项的模平方之和为 \(2n^2|z|^{2n-2}\)，极坐标积分给出
\[
 \int_D|\nabla U_n|^2\,\mathrm dx
 =2n^2\int_0^{2\pi}\int_0^1r^{2n-2}r\,\mathrm dr\,\mathrm d\theta
 =2\pi n.
\]
在这个显式函数族上，\(Q(g_n)\) 恰为内部梯度能量的 \(2\pi\) 倍，而边界 \(L^2\) 范数不随 \(n\) 增长。这反映了分数阶边界量对振荡的敏感性。本题计算的是指定的弦长双积分；正文通过固定图表和分割选取的边界范数，数值还取决于那些选择。
\end{solution}

\begin{exercise}\label{b1:ex:23-higher-zero-endpoints}
设 \(I=(0,1)\)、\(k\ge1\) 为整数、\(1\le p<\infty\)。以
\[
 W_0^{k,p}(I)=\overline{C_c^\infty(I)}^{\,W^{k,p}(I)}
\]
定义高阶零边界空间。给定 \(u\in W^{k,p}(I)\)，对 \(0\le j<k\)，用
\(u^{(j)}\) 的绝对连续代表定义它在 \(0,1\) 处的端点值。证明下列条件等价。
\begin{equivalences}
 \item\label{b1:item:23-higher-zero-closure}
 \(u\in W_0^{k,p}(I)\)。
 \item\label{b1:item:23-higher-zero-jets}
 对所有 \(0\le j<k\)，都有 \(u^{(j)}(0)=u^{(j)}(1)=0\)。
\end{equivalences}
证明充分性时，应给出边界截断的高阶导数估计；仅仅说“在端点附近截去一小段”不足以推出 \(W^{k,p}\) 收敛。
\end{exercise}
% 来源：自拟；一维绝对连续积分表示与内部紧支集逼近的高阶组合，不借助一般高阶迹空间。
\begin{solution}
每个 \(u^{(j)}\)、\(j<k\)，都属于 \(W^{1,p}(I)\)，因而具有题设的绝对连续代表。若 \(u_n\in C_c^\infty(I)\) 且 \(u_n\to u\) 于 \(W^{k,p}(I)\)，则
\[
 u_n^{(j)}\longrightarrow u^{(j)}
       \quad\text{于 }W^{1,p}(I),\quad 0\le j<k.
\]
一维代表估计，即命题\ref{b1:prop:sobolev-interval-representative}，使这个收敛也成为闭区间上的一致收敛。由于 \(u_n^{(j)}\) 的两端值都是零，极限的两端值也为零，必要性得证。

反过来，设所有这些端点值都为零。反复使用微积分基本定理，对于 \(0\le r<k\)，有
\[
 u^{(r)}(x)=\frac1{(k-r-1)!}
             \int_0^x(x-t)^{k-r-1}u^{(k)}(t)\,\mathrm dt.
\]
将右端视为截断积分核与 \(u^{(k)}\) 的卷积，并使用 Minkowski 积分不等式，可得对 \(0<a<1\)，
\[
 \|u^{(r)}\|_{L^p(0,a)}
 \le\frac{a^{k-r}}{(k-r)!}\,
             \|u^{(k)}\|_{L^p(0,a)}.
\]
具体地，只需把 \(u^{(k)}|_{(0,a)}\) 向外补零；核
\(t^{k-r-1}\mathbf1_{(0,a)}(t)/(k-r-1)!\) 的 \(L^1\) 范数为
\(a^{k-r}/(k-r)!\)，平移保持 \(L^p\) 范数，于是逐个平移积分便得到上式。这一步同样适用于 \(p=1\)。从右端点向左积分，还得到
\[
 \|u^{(r)}\|_{L^p(1-a,1)}
 \le\frac{a^{k-r}}{(k-r)!}\,
             \|u^{(k)}\|_{L^p(1-a,1)}.
\]

取 \(0<\varepsilon<1/4\)，选取 \(0\le\eta_\varepsilon\le1\) 的光滑截断，使它在
\([2\varepsilon,1-2\varepsilon]\) 上为 \(1\)，在两个端点的 \(\varepsilon\) 邻域内为零，并满足
\[
 \|\eta_\varepsilon^{(\ell)}\|_\infty
          \le C_\ell\varepsilon^{-\ell},\quad 1\le\ell\le k.
\]
记 \(A_\varepsilon=(0,2\varepsilon)\cup(1-2\varepsilon,1)\)。
对 \(0\le j\le k\)，Leibniz 公式给出
\[
 \bigl((1-\eta_\varepsilon)u\bigr)^{(j)}
 =(1-\eta_\varepsilon)u^{(j)}
  -\sum_{\ell=1}^j\binom j\ell
           \eta_\varepsilon^{(\ell)}u^{(j-\ell)}.
\]
第一项的 \(L^p\) 范数不超过
\(\|u^{(j)}\|_{L^p(A_\varepsilon)}\)，因 \(p<\infty\) 而趋于零。和式中
\(j-\ell<k\)，可以应用刚才的端点估计，得到
\[
 \begin{aligned}
 \|\eta_\varepsilon^{(\ell)}u^{(j-\ell)}\|_{L^p(I)}
 &\le C\varepsilon^{-\ell}
       \|u^{(j-\ell)}\|_{L^p(A_\varepsilon)}\\
 &\le C\varepsilon^{k-j}
       \|u^{(k)}\|_{L^p(A_\varepsilon)}
 \longrightarrow0.
 \end{aligned}
\]
即使 \(j=k\)，最后一项也因 \(u^{(k)}\in L^p(I)\) 的积分绝对连续性而趋于零。因此
\(\eta_\varepsilon u\to u\) 于 \(W^{k,p}(I)\)。

每个 \(\eta_\varepsilon u\) 都有包含在 \(I\) 内部的紧支集。由命题%
\ref{b1:prop:compact-sobolev-smooth-approximation}，它可以在 \(W^{k,p}(I)\) 中由
\(C_c^\infty(I)\) 函数逼近。先选 \(\varepsilon\) 控制截断误差，再作光滑逼近，便得到
\(u\in W_0^{k,p}(I)\)。这里有限 \(p\) 保证了边界层上的最高阶导数范数趋零；证明并未涵盖 \(p=\infty\)。
\end{solution}

\begin{exercise}\label{b1:ex:23-smooth-normal-data}
设 \(d\ge2\)，\(H=\mathbb R^{d-1}\times(0,\infty)\)，其边界外法向为
\(\nu=-e_d\)。本题只考虑光滑边界数据。
\begin{enumerate}
 \item\label{b1:item:23-independent-smooth-normal-data}
 给定任意 \(a,b\in C_c^\infty(\mathbb R^{d-1})\)，构造一个在
 \(\overline H\) 附近光滑、支集在 \(\overline H\) 中紧的函数 \(U\)，满足
 \[
  U(x',0)=a(x'),\quad \partial_\nu U(x',0)=b(x').
 \]
 说明该函数属于每个整数阶的 \(W^{k,p}(H)\)，其中 \(1\le p<\infty\)。
 \item\label{b1:item:23-smooth-tangential-compatibility}
 若还要求对 \(1\le i<d\) 有
 \(\partial_{x_i}U(x',0)=c_i(x')\)，其中 \(c_i\in C_c^\infty(\mathbb R^{d-1})\)，证明这些要求可同时满足，当且仅当
 \(c_i=\partial_{x_i}a\) 对每个 \(i\) 成立。
\end{enumerate}
\end{exercise}
% 来源：自拟；显式光滑边界数据延拓，分开自由法向数据与受约束的切向数据，不声称给出了粗糙数据的高阶右逆。
\begin{solution}
取 \(\chi\in C_c^\infty(\mathbb R)\)，使 \(\chi\) 在 \(0\) 的邻域恒等于 \(1\)。令
\[
 U(x',t)=\chi(t)\bigl(a(x')-t b(x')\bigr),\quad t\ge0.
\]
这个公式也在 \(t<0\) 定义光滑函数，因而确实给出了边界附近的光滑延拓。它在
\(\overline H\) 中的支集包含于
\[
 \bigl(\operatorname{supp}a\cup\operatorname{supp}b\bigr)
       \times\bigl(\operatorname{supp}\chi\cap[0,\infty)\bigr),
\]
后者为紧集。边界处 \(\chi(0)=1\)、\(\chi'(0)=0\)，所以
\[
 U(x',0)=a(x'),\quad
 \partial_tU(x',0)=-b(x'),\quad
 \partial_\nu U(x',0)=-\partial_tU(x',0)=b(x').
\]
所有阶导数都是紧支集的有界光滑函数，因此属于每个有限指数的 \(L^p(H)\)，也就得到题设的全部整数阶 Sobolev 正则性。

对任何在边界附近光滑的 \(U\)，把恒等式 \(U(x',0)=a(x')\) 沿边界变量求导，必有
\[
 \partial_{x_i}U(x',0)=\partial_{x_i}a(x'),\quad 1\le i<d.
\]
这证明条件的必要性。若 \(c_i=\partial_{x_i}a\)，第一问构造的 \(U\) 已经同时满足它们，故条件也充分。由此可见，光滑函数的边界值与法向导数可以独立指定，切向导数却已由边界值决定。整个构造只处理光滑数据，不据此推断一般分数阶数据所需的高阶延拓估计。
\end{solution}
